\chapter{Mean-field approaches}
\label{chap:hartrefock}
%=============================================================================

Chapter~\ref{chap:fci} showed that the exact answer is a superposition of an
exponentially large number of Slater determinants, and that we cannot afford
it.  This chapter takes the opposite tack.  Instead of adding determinants we
keep exactly one and ask which one is best.  The answer is Hartree-Fock
theory, and it is the foundation on which every method in the remaining
chapters is built: perturbation theory expands about it, coupled cluster
exponentiates corrections to it, and the truncated configuration interaction
of Section~\ref{sec:citruncations} is measured against it.

Chapter~\ref{chap:mbphys} already wrote down the Hartree-Fock equations,
Eq.~(\ref{eq:2-hfequations}), and deferred their derivation.  We supply it
here, in two independent ways -- by varying the expansion coefficients of the
orbitals, and by varying the determinant itself in second quantization.  The
second route requires Thouless' theorem, which characterises what a
``nearby determinant'' means, and Thouless' theorem in turn rests on the
Baker-Campbell-Hausdorff expansion.  That expansion is worth a section of its
own, because the same algebra reappears in the Trotter splitting used for time
evolution and for quantum simulation later in the book.  Armed with it we can
ask not merely whether the Hartree-Fock solution is stationary but whether it
is a minimum, and we close with the one many-body system for which the whole
programme can be carried through analytically: the infinite homogeneous
electron gas.

%----------------------------------------------------------------------
\section{Why a mean field?}
\label{sec:whyhf}

\index{Hartree-Fock theory}
Hartree-Fock theory is an algorithm for finding an approximate ground state of
a given Hamiltonian.  It has two ingredients.  The first is a single-particle
basis defined by its own eigenvalue problem,
\begin{equation}
  \hat{h}^{\mathrm{HF}}\psi_{\alpha} = \varepsilon_{\alpha}\psi_{\alpha},
  \qquad
  \hat{h}^{\mathrm{HF}} = \hat{t} + \hat{u}_{\mathrm{ext}}
    + \hat{u}^{\mathrm{HF}},
  \label{eq:6-hfhamiltonian}
\end{equation}
in which $\hat{u}^{\mathrm{HF}}$ is a one-body potential still to be
determined.  The second is the prescription for determining it: choose
$\hat{u}^{\mathrm{HF}}$ so that
\begin{equation}
  E^{\mathrm{HF}} = \bracket{\Phi_0}{\hat{H}}{\Phi_0}
  \label{eq:6-hfenergy}
\end{equation}
is stationary, with $\bket{\Phi_0}$ a single Slater determinant built from the
$N$ lowest solutions of Eq.~(\ref{eq:6-hfhamiltonian}).  The variational
principle of Section~\ref{sec:variational} then guarantees
$E^{\mathrm{HF}}\ge E_0$, so whatever we obtain is an upper bound on the exact
ground-state energy.

We shall see that the resulting one-body operator is precisely the operator
$\hat{f}$ that appears when the Hamiltonian is normal-ordered with respect to
the reference determinant.  For a specific matrix element,
\begin{equation}
  \bracket{p}{\hat{h}^{\mathrm{HF}}}{q}
  = \bracket{p}{\hat{f}}{q}
  = \bracket{p}{\hat{t}+\hat{u}_{\mathrm{ext}}}{q}
    + \sum_{i\le F}\langle pi\vert\hat{v}\vert qi\rangle_{\mathrm{AS}},
  \label{eq:6-fockmatrixelement}
\end{equation}
so that
\begin{equation}
  \bracket{p}{\hat{u}^{\mathrm{HF}}}{q}
  = \sum_{i\le F}\langle pi\vert\hat{v}\vert qi\rangle_{\mathrm{AS}} .
  \label{eq:6-hfpotential}
\end{equation}
The sum runs over all single-particle states below the Fermi level, and this
is the whole content of the mean-field idea: the two-body interaction, summed
over the other fermions, creates an effective one-body field in which a given
fermion moves.

Two features of Eq.~(\ref{eq:6-hfpotential}) deserve emphasis.  It brings an
explicit \emph{medium dependence}, since the potential felt by a particle
depends on which states are occupied; and it brings an explicit dependence on
the interaction, so different interactions give different mean fields even at
the same density.  For atoms and molecules there is an external potential
$\hat{u}_{\mathrm{ext}}$ -- the attraction of the nuclei.  For nuclei there is
none: a nucleus is a self-bound system, held together by the interaction
between its own constituents, and the entire one-body field must be generated
by Eq.~(\ref{eq:6-hfpotential}).

%----------------------------------------------------------------------
\section{Variational calculus and Lagrange multipliers}
\label{sec:varcalculus}

\index{variational calculus}\index{Lagrange multipliers}
Before deriving the Hartree-Fock equations we recall the machinery, in its
simplest setting.  The calculus of variations deals with problems in which the
quantity to be made stationary is an integral,
\begin{equation}
  E[\Phi] = \int_a^b f\!\left(\Phi(x),
    \frac{\partial\Phi}{\partial x}, x\right)\ud{x} .
  \label{eq:6-functional}
\end{equation}
The difficulty is that although $f$ is a known function of $\Phi$,
$\partial\Phi/\partial x$ and $x$, the dependence of $\Phi$ on $x$ is not
known: the limits are fixed but the path is not.  In our case the unknowns are
the single-particle wave functions, and the constraints are the integrals
expressing their orthonormality.  Constraints of this kind are handled with
Lagrange multipliers.

Take one particle in three dimensions, with
\begin{equation}
  E = \int \psi^{*}\hat{H}\psi\,\ud{x}\ud{y}\ud{z},
  \qquad
  \int \psi^{*}\psi\,\ud{x}\ud{y}\ud{z} = 1,
  \qquad
  \hat{H} = -\frac{1}{2}\nabla^2 + V .
  \label{eq:6-oneparticle}
\end{equation}
Integrating the kinetic term by parts, with $\psi$ vanishing at the
boundaries,
\begin{equation}
  \int \psi^{*}\left(-\frac{1}{2}\nabla^2\right)\psi\,\ud{x}\ud{y}\ud{z}
  = \int \frac{1}{2}\nabla\psi^{*}\cdot\nabla\psi\,\ud{x}\ud{y}\ud{z},
  \label{eq:6-byparts}
\end{equation}
so that the constrained variation reads
\begin{equation}
  \delta\left\{\int\left(\frac{1}{2}\nabla\psi^{*}\cdot\nabla\psi
    + V\psi^{*}\psi - \lambda\psi^{*}\psi\right)
    \ud{x}\ud{y}\ud{z}\right\} = 0 .
  \label{eq:6-constrainedvariation}
\end{equation}
Writing the integrand as
\begin{equation}
  f = \frac{1}{2}\left(\psi^{*}_x\psi_x + \psi^{*}_y\psi_y
      + \psi^{*}_z\psi_z\right) + V\psi^{*}\psi - \lambda\psi^{*}\psi,
  \label{eq:6-integrand}
\end{equation}
with subscripts denoting partial derivatives, the Euler-Lagrange equation for
$\psi^{*}$,
\begin{equation}
  \frac{\partial f}{\partial\psi^{*}}
  - \frac{\partial}{\partial x}\frac{\partial f}{\partial\psi^{*}_x}
  - \frac{\partial}{\partial y}\frac{\partial f}{\partial\psi^{*}_y}
  - \frac{\partial}{\partial z}\frac{\partial f}{\partial\psi^{*}_z} = 0,
  \label{eq:6-eulerlagrange}
\end{equation}
gives
\begin{equation}
  -\frac{1}{2}\nabla^2\psi + V\psi = \lambda\psi .
  \label{eq:6-schrodinger}
\end{equation}
The multiplier is the energy, and the stationarity condition is nothing but
the Schr\"odinger equation.  This is the pattern we shall meet again: a
constrained variation produces an eigenvalue problem, and the multiplier
enforcing normalisation turns out to be the eigenvalue.  It happened already
in Section~\ref{sec:cieigenvalue}, where the same manipulation produced the
full configuration-interaction equation, and it will happen once more below.

%----------------------------------------------------------------------
\section{Determinants under a linear change of basis}
\label{sec:detbasis}

\index{determinant!under a change of basis}
One property of determinants is needed before we can vary the coefficients,
and it is worth isolating because it explains why the whole construction is
consistent.

A determinant is linear in each column separately.  For a $2\times2$ example,
\begin{equation}
  \begin{vmatrix}
    \alpha_1 b_{11}+\alpha_2 b_{12} & a_{12}\\
    \alpha_1 b_{21}+\alpha_2 b_{22} & a_{22}
  \end{vmatrix}
  = \alpha_1\begin{vmatrix} b_{11} & a_{12}\\ b_{21} & a_{22}\end{vmatrix}
  + \alpha_2\begin{vmatrix} b_{12} & a_{12}\\ b_{22} & a_{22}\end{vmatrix},
  \label{eq:6-multilinear2}
\end{equation}
and in general, for an $n\times n$ matrix in which one column is a linear
combination $\sum_k c_k b_{jk}$,
\begin{equation}
  \det\left(\dots\ \textstyle\sum_k c_k b_{jk}\ \dots\right)
  = \sum_{k=1}^{n} c_k \det\left(\dots\ b_{jk}\ \dots\right),
  \label{eq:6-multilinear}
\end{equation}
where the omitted columns are the same on both sides.  Applying this to
\emph{every} column in turn gives the product rule
\begin{equation}
  \det(\bm{B}\bm{C}) = \det(\bm{B})\,\det(\bm{C}),
  \label{eq:6-productrule}
\end{equation}
which we already used in Section~\ref{sec:determinants}.  The reader is
encouraged to check the case $n=2$ explicitly.

Now let the new orbitals be expanded in a fixed orthonormal basis, as in
Eq.~(\ref{eq:newbasis}),
\begin{equation}
  \psi_p(x) = \sum_{\lambda}C_{p\lambda}\phi_{\lambda}(x).
  \label{eq:6-newbasis}
\end{equation}
The Slater determinant built from the $\psi$'s is
\begin{equation}
  \frac{1}{\sqrt{N!}}
  \begin{vmatrix}
    \sum_{\lambda}C_{p\lambda}\phi_{\lambda}(x_1) & \dots
      & \sum_{\lambda}C_{p\lambda}\phi_{\lambda}(x_N)\\
    \vdots & & \vdots\\
    \sum_{\lambda}C_{t\lambda}\phi_{\lambda}(x_1) & \dots
      & \sum_{\lambda}C_{t\lambda}\phi_{\lambda}(x_N)
  \end{vmatrix}
  = \det(\bm{C})\;\Phi(\phi),
  \label{eq:6-detcdetphi}
\end{equation}
with $\Phi(\phi)$ the determinant of the original basis functions.  If
$\bm{C}$ is unitary then $|\det(\bm{C})|=1$ and the two determinants differ by
a phase only.  Two consequences follow, and both are used constantly.

First, orthonormality is preserved.  If $\bm{v}_j^{T}\bm{v}_i=\delta_{ij}$ and
$\bm{w}_i=\bm{U}\bm{v}_i$ with $\bm{U}$ unitary, then
\begin{equation}
  \bm{w}_j^{\dagger}\bm{w}_i
  = \bm{v}_j^{\dagger}\bm{U}^{\dagger}\bm{U}\bm{v}_i
  = \bm{v}_j^{\dagger}\bm{v}_i = \delta_{ij},
  \label{eq:6-unitarypreserves}
\end{equation}
which is Section~\ref{sec:unitary} restated.  So if $\braket{\alpha}{\beta}
=\delta_{\alpha\beta}$ then $\braket{p}{q}=\delta_{pq}$, and a unitary change
of coefficients never spoils the basis we are working in.

Second, a Slater determinant is invariant, up to a phase, under any unitary
mixing of its occupied orbitals.  The determinant does not know which
particular orbitals were used to build it, only which $N$-dimensional subspace
they span.  This is why the Hartree-Fock solution is properly a statement
about a subspace, and why the individual orbitals are determined only up to
such a mixing -- the freedom that is fixed, conventionally, by demanding that
the Fock matrix be diagonal.

\begin{notebox}
\textbf{Where this comes from.}  Every ingredient of this section already
appeared in Chapter~\ref{chap:linalg}: multilinearity and the product rule
in Section~\ref{sec:determinants}, unitary transformations in
Section~\ref{sec:unitary}, and the practical evaluation of
$\det(\bm{C})$ by LU decomposition in Section~\ref{sec:lu}.  The last is not
an idle remark: in variational Monte Carlo the ratio of two such determinants
is evaluated millions of times, which is why
Section~\ref{sec:slaterupdate} developed the Sherman-Morrison update instead.
\end{notebox}

%----------------------------------------------------------------------
\section{The Hartree-Fock equations by varying the coefficients}
\label{sec:hfcoefficients}

\index{Hartree-Fock equations!derivation}
We can now carry out the derivation that Chapter~\ref{chap:mbphys} postponed.
The energy functional of a single determinant, Eq.~(\ref{eq:EHF}), is
\begin{equation}
  E[\Phi] = \sum_{\mu=1}^{N}\bracket{\mu}{\hat{h}_0}{\mu}
    + \frac{1}{2}\sum_{\mu=1}^{N}\sum_{\nu=1}^{N}
      \langle\mu\nu\vert\hat{v}\vert\mu\nu\rangle_{\mathrm{AS}} ,
  \label{eq:6-energyfunctional}
\end{equation}
and inserting Eq.~(\ref{eq:6-newbasis}) turns it into a function of the
coefficients,
\begin{equation}
  E[\Phi^{\mathrm{HF}}]
  = \sum_{i=1}^{N}\sum_{\alpha\beta}C^{*}_{i\alpha}C_{i\beta}
      \bracket{\alpha}{\hat{h}_0}{\beta}
  + \frac{1}{2}\sum_{i,j=1}^{N}\sum_{\alpha\beta\gamma\delta}
      C^{*}_{i\alpha}C^{*}_{j\beta}C_{i\gamma}C_{j\delta}
      \langle\alpha\beta\vert\hat{v}\vert\gamma\delta\rangle_{\mathrm{AS}} .
  \label{eq:6-Efunctional}
\end{equation}
Greek indices run over the whole fixed basis, Latin indices over the $N$
occupied orbitals.

The coefficients are not free: orthonormality of the new orbitals requires
\begin{equation}
  \braket{i}{j} = \delta_{ij}
  = \sum_{\alpha\beta}C^{*}_{i\alpha}C_{i\beta}\braket{\alpha}{\beta}
  = \sum_{\alpha}C^{*}_{i\alpha}C_{i\alpha},
  \label{eq:6-cconstraint}
\end{equation}
so we introduce one multiplier $\epsilon_i$ per occupied orbital and minimise
\begin{equation}
  F[\Phi^{\mathrm{HF}}] = E[\Phi^{\mathrm{HF}}]
    - \sum_{i=1}^{N}\epsilon_i\sum_{\alpha}C^{*}_{i\alpha}C_{i\alpha}.
  \label{eq:6-lagrangian}
\end{equation}
The coefficients and their complex conjugates may be varied independently --
the real and imaginary parts are two independent real variations, and it is
equivalent and more convenient to use $C$ and $C^{*}$ -- so we impose
\begin{equation}
  \frac{\partial}{\partial C^{*}_{i\alpha}}
  \left[E[\Phi^{\mathrm{HF}}]
    - \sum_{j}\epsilon_j\sum_{\beta}C^{*}_{j\beta}C_{j\beta}\right] = 0 .
  \label{eq:6-variationC}
\end{equation}
Differentiating Eq.~(\ref{eq:6-Efunctional}) is straightforward.  The one-body
term contributes $\sum_{\beta}C_{i\beta}\bracket{\alpha}{\hat{h}_0}{\beta}$.
In the two-body term $C^{*}_{i\alpha}$ appears twice, once as
$C^{*}_{i\alpha}$ and once as $C^{*}_{j\beta}$ with $j=i$; the two
contributions are equal by the antisymmetry of
$\langle\alpha\beta\vert\hat{v}\vert\gamma\delta\rangle_{\mathrm{AS}}$ under
the simultaneous interchange of both pairs, and they cancel the factor
$\tfrac12$.  The result is
\begin{equation}
  \sum_{\beta}C_{i\beta}\bracket{\alpha}{\hat{h}_0}{\beta}
  + \sum_{j=1}^{N}\sum_{\beta\gamma\delta}
      C^{*}_{j\beta}C_{j\delta}C_{i\gamma}
      \langle\alpha\beta\vert\hat{v}\vert\gamma\delta\rangle_{\mathrm{AS}}
  = \epsilon_i^{\mathrm{HF}}C_{i\alpha}.
  \label{eq:6-hfraw}
\end{equation}
Relabelling the dummy indices so that the coefficient $C_{i\beta}$ can be
factored out of both terms,
\begin{equation}
  \sum_{\beta}\left\{\bracket{\alpha}{\hat{h}_0}{\beta}
    + \sum_{j=1}^{N}\sum_{\gamma\delta}C^{*}_{j\gamma}C_{j\delta}
      \langle\alpha\gamma\vert\hat{v}\vert\beta\delta\rangle_{\mathrm{AS}}
    \right\}C_{i\beta} = \epsilon_i^{\mathrm{HF}}C_{i\alpha}.
  \label{eq:6-hfbraced}
\end{equation}
Defining the Hartree-Fock, or Fock, matrix
\begin{equation}
  h^{\mathrm{HF}}_{\alpha\beta}
  = \bracket{\alpha}{\hat{h}_0}{\beta}
  + \sum_{j=1}^{N}\sum_{\gamma\delta}C^{*}_{j\gamma}C_{j\delta}
    \langle\alpha\gamma\vert\hat{v}\vert\beta\delta\rangle_{\mathrm{AS}},
  \label{eq:6-fockmatrix}
\end{equation}
the condition becomes
\begin{equation}
  \boxed{\;
  \sum_{\beta}h^{\mathrm{HF}}_{\alpha\beta}C_{i\beta}
    = \epsilon_i^{\mathrm{HF}}C_{i\alpha}
  \;}
  \label{eq:6-hfequations}
\end{equation}
This is Eq.~(\ref{eq:2-hfequations}), now derived.  It is a standard Hermitian
eigenvalue problem, solvable by the Householder reduction and QL iteration of
Sections~\ref{sec:householder} and~\ref{sec:qlalgorithm}, or by Lanczos,
Section~\ref{sec:lanczos}, when the basis is large -- with the one
complication that the matrix depends on its own eigenvectors.

If the fixed basis is chosen as the eigenbasis of $\hat{h}_0$, which is the
usual choice, the first term is diagonal,
\begin{equation}
  h^{\mathrm{HF}}_{\alpha\beta}
  = \epsilon_{\alpha}\delta_{\alpha\beta}
  + \sum_{j=1}^{N}\sum_{\gamma\delta}C^{*}_{j\gamma}C_{j\delta}
    \langle\alpha\gamma\vert\hat{v}\vert\beta\delta\rangle_{\mathrm{AS}},
  \label{eq:6-fockdiagonal}
\end{equation}
and Eq.~(\ref{eq:6-hfequations}) can be written compactly as
$\bm{h}^{\mathrm{HF}}\bm{C} = \bm{\epsilon}^{\mathrm{HF}}\bm{C}$.

An important practical point is visible in Eq.~(\ref{eq:6-fockdiagonal}): no
integrals need to be recomputed during the iteration.  The matrix elements
$\bracket{\alpha}{\hat{h}_0}{\beta}$ and
$\langle\alpha\gamma\vert\hat{v}\vert\beta\delta\rangle_{\mathrm{AS}}$ are
tabulated once and for all; every subsequent iteration is a contraction of
those tables with the current coefficients, followed by a diagonalisation.
The cost of storing them is the $n^4$ scaling discussed in
Section~\ref{sec:twobodycost}, and it is the reason the factorisations of
Section~\ref{sec:twobodyfact} matter.

%----------------------------------------------------------------------
\section{The density matrix and the self-consistent field algorithm}
\label{sec:scf}

\index{density matrix}\index{self-consistent field}
The double sum over the occupied orbitals in Eq.~(\ref{eq:6-fockdiagonal})
depends on the coefficients only through the combination
\begin{equation}
  \rho_{\gamma\delta}
  = \sum_{i=1}^{N}\braket{\gamma}{i}\braket{i}{\delta}
  = \sum_{i=1}^{N}C_{i\gamma}C^{*}_{i\delta},
  \label{eq:6-densitymatrix}
\end{equation}
the one-body \emph{density matrix}, which is precisely the object introduced
in Section~\ref{sec:naturalorbitals}.  It is the projector onto the occupied
subspace, $\bm{\rho}^2=\bm{\rho}$ and $\mathrm{tr}\,\bm{\rho}=N$, which is the
formal statement of the invariance found in Section~\ref{sec:detbasis}: the
Fock matrix depends on the occupied subspace, not on the particular orbitals
chosen to span it.  With it,
\begin{equation}
  h^{\mathrm{HF}}_{\alpha\beta}
  = \epsilon_{\alpha}\delta_{\alpha\beta}
  + \sum_{\gamma\delta}\rho_{\gamma\delta}
    \langle\alpha\gamma\vert\hat{v}\vert\beta\delta\rangle_{\mathrm{AS}} .
  \label{eq:6-fockdensity}
\end{equation}
Building $\bm{\rho}$ once per iteration and contracting it with the stored
two-body table is the efficient way to organise the calculation.

The algorithm is then an iteration of the kind discussed in
Section~\ref{sec:iterative}, with a diagonalisation replacing a single sweep.

\begin{enumerate}
  \item Start from a guess, conventionally $C^{(0)}_{i\alpha}
    =\delta_{i\alpha}$, which makes the reference determinant the one built
    from the $N$ lowest eigenstates of $\hat{h}_0$.  Random orthonormal
    vectors work too, and are sometimes preferable because they avoid
    convergence onto a symmetry-constrained saddle point.
  \item Build $\bm{\rho}$ from Eq.~(\ref{eq:6-densitymatrix}) and
    $\bm{h}^{\mathrm{HF}}$ from Eq.~(\ref{eq:6-fockdensity}).
  \item Diagonalise, obtaining $C^{(1)}_{i\alpha}$ and
    $\epsilon_i^{\mathrm{HF}(1)}$.
  \item Repeat from step 2 until the single-particle energies stop moving,
    \begin{equation}
      \frac{1}{m}\sum_{p}\left\vert
        \epsilon_p^{(n)}-\epsilon_p^{(n-1)}\right\vert \le \lambda,
      \label{eq:6-convergence}
    \end{equation}
    with $m$ the number of single-particle states and $\lambda$ a prescribed
    tolerance, typically $10^{-8}$ or smaller.
\end{enumerate}

The program \texttt{hartreefock.py} implements exactly this loop.  For the
system of Section~\ref{sec:energyfunctional} -- spinless fermions in a
one-dimensional oscillator trap with a softened Coulomb repulsion, in a basis
of eight oscillator orbitals -- it reproduces Table~\ref{tab:hfgain}, and the
gain over the unoptimised determinant grows with the particle number:
\begin{center}
\renewcommand{\arraystretch}{1.2}
\begin{tabular}{rrlll}
\toprule
$N$ & iterations & $E$(reference) & $E$(Hartree-Fock) & gain\\
\midrule
$2$ & $17$ & $2.69018342$  & $2.66308878$  & $-2.71\times10^{-2}$\\
$3$ & $17$ & $6.46801021$  & $6.39016133$  & $-7.78\times10^{-2}$\\
$4$ & $17$ & $11.77084712$ & $11.62305385$ & $-1.48\times10^{-1}$\\
\bottomrule
\end{tabular}
\end{center}

\paragraph{Brillouin's theorem.}
\index{Brillouin's theorem}
At convergence the Fock matrix is diagonal in its own eigenbasis, so its
occupied-unoccupied block vanishes,
\begin{equation}
  f_{ai} = 0
  \qquad\Longleftrightarrow\qquad
  \bracket{\Phi_0}{\hat{H}}{\Phi_i^a} = 0 .
  \label{eq:6-brillouin}
\end{equation}
The program reports $\max|f_{ai}| = 8.7\times10^{-11}$ for $N=4$, at the
prescribed tolerance.  This is Brillouin's theorem, met first in
Section~\ref{sec:wickapps} and used in Section~\ref{sec:cimatrix} to empty the
$0p0h$--$1p1h$ block of the configuration-interaction matrix.  We can now see
what it means: the Hartree-Fock solution is the choice of single-particle
basis that decouples the reference from all single excitations.  Everything
left over is carried by the doubles and higher, which is why
Eq.~(\ref{eq:5-energyeq}) contained only doubles and why both perturbation
theory and coupled-cluster theory start there.

%----------------------------------------------------------------------
\section{Hartree-Fock in second quantization}
\label{sec:hfsecondquant}

The derivation of Section~\ref{sec:hfcoefficients} varied the orbitals.  We
now vary the determinant, which is more abstract but far better suited to the
question of stability, and which is the natural language given
Chapter~\ref{chap:secondquant}.

Our ansatz is the determinant
\begin{equation}
  \bket{\Phi_0} = \bket{c}
  = \hat{a}^{\dagger}_{i}\hat{a}^{\dagger}_{j}\cdots
    \hat{a}^{\dagger}_{l}\bket{0},
  \label{eq:6-cstate}
\end{equation}
and we seek $\hat{u}^{\mathrm{HF}}$ such that
$E_0^{\mathrm{HF}}=\bracket{c}{\hat{H}}{c}$ is a local minimum.  The
stationarity condition is that the energy does not change to first order under
any variation of the determinant,
\begin{equation}
  \bracket{\delta c}{\hat{H}}{c} = 0
  \label{eq:6-stationarity}
\end{equation}
for every admissible $\bket{\delta c}$.  To make this precise we must know
what the admissible variations are -- that is, what the neighbourhood of a
Slater determinant looks like inside the space of Slater determinants.  That
is the content of Thouless' theorem.

%----------------------------------------------------------------------
\section{Thouless' theorem}
\label{sec:thouless}

\index{Thouless' theorem}
\begin{theorem}[Thouless]
Any Slater determinant $\bket{c'}$ that is not orthogonal to a given
determinant $\bket{c}=\prod_{i=1}^{n}\hat{a}^{\dagger}_{\alpha_i}\bket{0}$ can
be written as
\begin{equation}
  \bket{c'} = \exp\left\{\sum_{a>F}\sum_{i\le F}
    C_{ai}\,\hat{a}^{\dagger}_{a}\hat{a}_{i}\right\}\bket{c} .
  \label{eq:6-thouless}
\end{equation}
\end{theorem}

In words: every determinant in the neighbourhood of $\bket{c}$ is obtained
from it by an exponentiated linear combination of one-particle-one-hole
excitations.  This is a strong statement.  It says that the manifold of Slater
determinants is parametrised, near any point, by the $(n-N)\times N$
amplitudes $C_{ai}$ and by nothing else -- so the Hartree-Fock variation has
exactly that many degrees of freedom.

\paragraph{Step one: the exponential collapses.}
The exponential of a sum of operators factorises into a product of
exponentials only when the operators commute.  When they do not, the
Baker-Campbell-Hausdorff formula of Section~\ref{sec:bch} supplies the
correction.  Here we are lucky: for particle indices $a,b>F$ and hole indices
$i,j\le F$, the two index sets are disjoint, and a short calculation with the
anticommutation relations of Chapter~\ref{chap:secondquant} gives
\begin{equation}
  \left[\hat{a}^{\dagger}_{a}\hat{a}_{i},\,
        \hat{a}^{\dagger}_{b}\hat{a}_{j}\right] = 0 .
  \label{eq:6-phcommute}
\end{equation}
Indeed, since $i\ne b$ we may write
$\hat{a}_{i}\hat{a}^{\dagger}_{b}=-\hat{a}^{\dagger}_{b}\hat{a}_{i}$, and
likewise $\hat{a}_{j}\hat{a}^{\dagger}_{a}
=-\hat{a}^{\dagger}_{a}\hat{a}_{j}$; both orderings then reduce to
$-\hat{a}^{\dagger}_{a}\hat{a}^{\dagger}_{b}\hat{a}_{i}\hat{a}_{j}$ and the
commutator vanishes.  Writing
\begin{equation}
  \hat{T} = \sum_{i\le F}\hat{T}_i,
  \qquad
  \hat{T}_i = \sum_{a>F}C_{ai}\hat{a}^{\dagger}_{a}\hat{a}_{i},
  \label{eq:6-Tsplit}
\end{equation}
we therefore have $e^{\hat{T}}=\prod_i e^{\hat{T}_i}$.  Moreover each factor
truncates, because
\begin{equation}
  \hat{T}_i^2 = \sum_{ab>F}C_{ai}C_{bi}\,
    \hat{a}^{\dagger}_{a}\hat{a}_{i}\hat{a}^{\dagger}_{b}\hat{a}_{i} = 0 ,
  \label{eq:6-Tsquared}
\end{equation}
since the rightmost $\hat{a}_i$ empties the hole $i$ and the second one then
annihilates the state -- an instance of $(\hat{a}_i)^n=0$ for $n>1$.  Hence
\begin{equation}
  \bket{c'} = \prod_{i\le F}
    \left\{1+\sum_{a>F}C_{ai}\hat{a}^{\dagger}_{a}\hat{a}_{i}\right\}\bket{c}.
  \label{eq:6-thoulessproduct}
\end{equation}
Note carefully that the product in
Eq.~(\ref{eq:6-thoulessproduct}) is \emph{not} the same as
$1+\hat{T}$: the cross terms $\hat{T}_i\hat{T}_j$ with $i\ne j$ survive and
describe $2p$-$2h$ admixtures.  Dropping them is legitimate only for
infinitesimal amplitudes, which is exactly the situation in the stability
analysis below.

\paragraph{Step two: the product is a determinant.}
Insert $\bket{c}=\hat{a}^{\dagger}_{i_1}\hat{a}^{\dagger}_{i_2}\cdots
\hat{a}^{\dagger}_{i_n}\bket{0}$ into
Eq.~(\ref{eq:6-thoulessproduct}).  The factor labelled $i_1$ acts only on
$\hat{a}^{\dagger}_{i_1}$, the factor labelled $i_2$ only on
$\hat{a}^{\dagger}_{i_2}$, and so on, so the product reorganises into
\begin{align}
  \bket{c'}
  &= \left(1+\sum_{a>F}C_{ai_1}\hat{a}^{\dagger}_{a}\hat{a}_{i_1}\right)
     \hat{a}^{\dagger}_{i_1}
     \left(1+\sum_{a>F}C_{ai_2}\hat{a}^{\dagger}_{a}\hat{a}_{i_2}\right)
     \hat{a}^{\dagger}_{i_2}\cdots\bket{0}
     \nonumber\\
  &= \prod_{i\le F}\left(\hat{a}^{\dagger}_{i}
     + \sum_{a>F}C_{ai}\hat{a}^{\dagger}_{a}\right)\bket{0} .
  \label{eq:6-thoulessdet}
\end{align}
Defining a new creation operator
\begin{equation}
  \hat{b}^{\dagger}_{i} = \hat{a}^{\dagger}_{i}
    + \sum_{a>F}C_{ai}\hat{a}^{\dagger}_{a},
  \label{eq:6-newb}
\end{equation}
we have $\bket{c'}=\prod_i \hat{b}^{\dagger}_{i}\bket{0}$, which has exactly
the form of a Slater determinant.  The rotated state is a single determinant,
not a superposition: this is the whole point.

\paragraph{Step three: it is the general determinant.}
Equation~(\ref{eq:6-newb}) looks less general than it should, because the sum
runs only over states above the Fermi level.  Consider therefore the most
general one,
\begin{equation}
  \tilde{b}^{\dagger}_{i} = \sum_{p}f_{ip}\hat{a}^{\dagger}_{p},
  \qquad
  \bket{\tilde{c}} = \prod_{i}\tilde{b}^{\dagger}_{i}\bket{0},
  \label{eq:6-generalb}
\end{equation}
with $f$ unitary, and require only that
$\braket{c}{\tilde{c}}\ne0$.  Expanding,
\begin{equation}
  \braket{c}{\tilde{c}}
  = \bbra{0}\hat{a}_{i_n}\cdots\hat{a}_{i_1}
    \left(\sum_{p}f_{i_1p}\hat{a}^{\dagger}_{p}\right)\cdots
    \left(\sum_{t}f_{i_nt}\hat{a}^{\dagger}_{t}\right)\bket{0}
  = \det\left(f_{ip}\right)_{i,p\le F},
  \label{eq:6-overlapdet}
\end{equation}
the determinant of the occupied block of $f$: only the terms with $p$ below
the Fermi level survive the vacuum expectation value, and summing over the
permutations reconstructs the determinant.  The hypothesis
$\braket{c}{\tilde{c}}\ne0$ says exactly that this determinant is non-zero, so
the occupied block is invertible, and intermediate normalisation lets us
rescale to $\det(f_{ip})=1$.  With $f^{-1}$ available we can form
\begin{equation}
  \sum_{i}f^{-1}_{ki}\tilde{b}^{\dagger}_{i}
  = \sum_{i}\sum_{p}f^{-1}_{ki}f_{ip}\hat{a}^{\dagger}_{p}
  = \hat{a}^{\dagger}_{k}
    + \sum_{p>F}\left(\sum_{i\le F}f^{-1}_{ki}f_{ip}\right)
      \hat{a}^{\dagger}_{p},
  \label{eq:6-finverse}
\end{equation}
where the terms with $p\le F$ have collapsed to $\delta_{kp}$.  Setting
$C_{pk}=\sum_{i\le F}f^{-1}_{ki}f_{ip}$ returns
Eq.~(\ref{eq:6-newb}) exactly.  Since a unitary mixing of the
$\tilde{b}^{\dagger}$ changes the determinant only by a factor, as
Section~\ref{sec:detbasis} showed, $\bket{\tilde{c}}$ and $\bket{c'}$ are the
same state.  The theorem is proved.

\paragraph{Numerical check.}
The program \texttt{hartreefock.py} verifies each step explicitly for six
orbitals and three particles, with random amplitudes of size $0.3$:
$\max|[\hat{T}_i,\hat{T}_j]| = 2.8\times10^{-17}$,
$\max|\hat{T}_i^2| = 0$,
$\bket{c'}$ agrees with $\prod_i\hat{b}^{\dagger}_i\bket{0}$ to
$5.6\times10^{-17}$, the overlap $\braket{c}{c'}$ is exactly one, and
$\det(f_{ip})=1$.  The difference between $e^{\hat{T}}\bket{c}$ and
$(1+\hat{T})\bket{c}$ falls by a factor of four each time the amplitudes are
halved, confirming that the discarded piece is second order.

\paragraph{What it buys us.}
Thouless' theorem lets us write any nearby determinant as
\begin{equation}
  \bket{c'} = \bket{c} + \bket{\delta c},
  \qquad
  \bket{c'} \simeq \left(1+\sum_{ai}\delta C_{ai}
    \hat{a}^{\dagger}_{a}\hat{a}_{i}\right)\bket{c}
  \label{eq:6-firstvariation}
\end{equation}
to first order in the amplitudes.  The stationarity
condition~(\ref{eq:6-stationarity}) then reads
\begin{equation}
  \sum_{ai}\delta C^{*}_{ai}\,
    \bracket{c}{\hat{a}^{\dagger}_{i}\hat{a}_{a}\hat{H}}{c} = 0
  \quad\text{for arbitrary }\delta C
  \qquad\Longrightarrow\qquad
  \bracket{\Phi_i^a}{\hat{H}}{\Phi_0} = 0 ,
  \label{eq:6-hfcondition}
\end{equation}
which is Brillouin's theorem again, Eq.~(\ref{eq:6-brillouin}).  The two
derivations therefore agree: varying the orbitals and varying the determinant
give the same equations.  What the second route adds is the ability to go to
second order, which is the subject of Section~\ref{sec:hfstability}.

%----------------------------------------------------------------------
\section{The Baker-Campbell-Hausdorff expansion}
\label{sec:bch}

\index{Baker-Campbell-Hausdorff formula}
Thouless' theorem needed only the trivial case of the following result, but
the general case is used repeatedly later -- in the similarity transformation
of coupled-cluster theory, in the interaction picture, and in the Trotter
splitting of the next section -- so we develop it here.

For two operators $X$ and $Y$ that do not commute, $e^{X}e^{Y}\ne e^{X+Y}$.
The Baker-Campbell-Hausdorff formula states that the product is nevertheless
a single exponential, $e^{X}e^{Y}=e^{Z}$, with
\begin{equation}
  Z = X + Y + \frac{1}{2}[X,Y]
    + \frac{1}{12}\Big([X,[X,Y]] + [Y,[Y,X]]\Big)
    - \frac{1}{24}\big[Y,[X,[X,Y]]\big] + \cdots
  \label{eq:6-bch}
\end{equation}
Every term beyond the first is a nested commutator of $X$ and $Y$.  No
ordinary products appear, which is the essential structural statement: if $X$
and $Y$ belong to a Lie algebra, so does $Z$.  The coefficients
$\tfrac12,\tfrac1{12},\tfrac1{24},\dots$ involve Bernoulli numbers and were
first given in general by Dynkin.

\paragraph{Derivation of the first terms.}
Expand both sides in powers and match order by order.  With
$Z=X+Y+A_2+A_3+\cdots$, where $A_n$ collects the terms of total degree $n$,
\begin{equation}
  e^{X}e^{Y} = I + (X+Y)
    + \frac{1}{2}\left(X^2+2XY+Y^2\right)
    + \bigO(3),
  \qquad
  e^{Z} = I + Z + \frac{1}{2}Z^2 + \bigO(3).
  \label{eq:6-bchmatch}
\end{equation}
At first order, $A_1$ is absent and $Z^{(1)}=X+Y$.  At second order the
left-hand side supplies $\tfrac12(X^2+2XY+Y^2)$ while
$\tfrac12(X+Y)^2=\tfrac12(X^2+XY+YX+Y^2)$; the difference is
$\tfrac12(XY-YX)=\tfrac12[X,Y]$, which must therefore be supplied by $A_2$,
\begin{equation}
  A_2 = \frac{1}{2}[X,Y].
  \label{eq:6-bchA2}
\end{equation}
The same procedure at third order, using the Jacobi identity
$[X,[Y,W]]+[Y,[W,X]]+[W,[X,Y]]=0$ to reduce the independent structures, gives
\begin{equation}
  A_3 = \frac{1}{12}\big[X,[X,Y]\big] - \frac{1}{12}\big[Y,[X,Y]\big]
      = \frac{1}{12}\Big([X,[X,Y]] + [Y,[Y,X]]\Big).
  \label{eq:6-bchA3}
\end{equation}

\paragraph{When the series terminates.}
If $[X,Y]$ commutes with both $X$ and $Y$ -- in particular if it is a multiple
of the identity, $[X,Y]=c\,I$ -- then every higher commutator vanishes and
\begin{equation}
  e^{X}e^{Y} = \exp\left(X+Y+\tfrac12[X,Y]\right)
  \label{eq:6-bchcentral}
\end{equation}
is exact.  This is the case for canonically conjugate operators,
$[\hat{x},\hat{p}]=\imath\hbar$, from which the Weyl form of the
displacement operator follows, and for the ladder operators of the harmonic
oscillator, $[\hat{a},\hat{a}^{\dagger}]=1$.  It is also, trivially, the case
in Thouless' theorem, where by Eq.~(\ref{eq:6-phcommute}) the commutator is
zero and the series stops at the first term.

\paragraph{How good is a truncation?}
For the two nilpotent matrices
$X=\left(\begin{smallmatrix}0 & 0.1\\ 0 & 0\end{smallmatrix}\right)$ and
$Y=\left(\begin{smallmatrix}0 & 0\\ 0.1 & 0\end{smallmatrix}\right)$,
\texttt{hartreefock.py} gives
\begin{center}
\renewcommand{\arraystretch}{1.2}
\begin{tabular}{clr}
\toprule
order & terms kept & $\Vert e^{X}e^{Y}-e^{Z_n}\Vert$\\
\midrule
$1$ & $X+Y$ & $7.071\times10^{-3}$\\
$2$ & $+\ \tfrac12[X,Y]$ & $2.379\times10^{-4}$\\
$3$ & $+\ \tfrac1{12}\big([X,[X,Y]]+[Y,[Y,X]]\big)$ & $1.190\times10^{-5}$\\
$4$ & $-\ \tfrac1{24}\big[Y,[X,[X,Y]]\big]$ & $4.757\times10^{-7}$\\
\bottomrule
\end{tabular}
\end{center}
Each additional commutator gains between one and two orders of magnitude for
amplitudes of this size.  The series is asymptotic rather than convergent in
general, and the gain per order shrinks as the operators grow; the practical
lesson is that BCH is a small-parameter expansion, and the small parameter is
the size of the commutator relative to the operators themselves.

%----------------------------------------------------------------------
\section{The Trotter-Suzuki splitting}
\label{sec:trotter}

\index{Trotter formula}\index{Suzuki-Trotter decomposition}
The most common use of Eq.~(\ref{eq:6-bch}) is in reverse.  Suppose we can
exponentiate $H_1$ and $H_2$ separately but not their sum -- the standard
situation, with $H_1$ the kinetic energy, diagonal in momentum space, and
$H_2$ the potential, diagonal in coordinate space.  Setting
$X=-\imath H_1\Delta$ and $Y=-\imath H_2\Delta$ in
Eq.~(\ref{eq:6-bch}),
\begin{equation}
  e^{-\imath H_1\Delta}e^{-\imath H_2\Delta}
  = \exp\left(-\imath (H_1+H_2)\Delta
      - \frac{1}{2}[H_1,H_2]\Delta^2 + \bigO(\Delta^3)\right),
  \label{eq:6-trotterbch}
\end{equation}
so a single split step is correct to first order in $\Delta$, with an error
governed by the commutator.  Dividing the evolution into $N$ steps of length
$t/N$ gives the Lie-Trotter product formula
\begin{equation}
  e^{-\imath(H_1+H_2)t}
  = \lim_{N\to\infty}\left(e^{-\imath H_1 t/N}
      e^{-\imath H_2 t/N}\right)^{N},
  \label{eq:6-trotter}
\end{equation}
whose error at finite $N$ is $\bigO(t^2/N)$.

Symmetrising the step cancels the leading commutator.  The second-order, or
Strang, splitting
\begin{equation}
  S_2(\Delta) = e^{-\imath H_1\Delta/2}\,e^{-\imath H_2\Delta}\,
                e^{-\imath H_1\Delta/2}
  = e^{-\imath(H_1+H_2)\Delta} + \bigO(\Delta^3)
  \label{eq:6-strang}
\end{equation}
has a local error governed by the double commutators $[H_1,[H_1,H_2]]$ and
$[H_2,[H_1,H_2]]$, hence a global error $\bigO(t^3/N^2)$.  Suzuki's recursive
construction generalises this to any even order, at the cost of more
exponentials per step and, from fourth order on, of steps with negative
coefficients.

For a single qubit with $\hat{H}=\sigma_x+\sigma_z$, so that
$[\sigma_x,\sigma_z]=-2\imath\sigma_y\ne0$, and $t=1$,
\texttt{hartreefock.py} finds
\begin{center}
\renewcommand{\arraystretch}{1.2}
\begin{tabular}{rllll}
\toprule
$N$ & first order & ratio & second order & ratio\\
\midrule
$1$  & $1.130\times10^{0}$  & --- & $4.436\times10^{-1}$ & ---\\
$2$  & $5.125\times10^{-1}$ & $2.21$ & $1.006\times10^{-1}$ & $4.41$\\
$4$  & $2.493\times10^{-1}$ & $2.06$ & $2.443\times10^{-2}$ & $4.12$\\
$8$  & $1.238\times10^{-1}$ & $2.01$ & $6.061\times10^{-3}$ & $4.03$\\
$16$ & $6.177\times10^{-2}$ & $2.00$ & $1.513\times10^{-3}$ & $4.01$\\
$32$ & $3.087\times10^{-2}$ & $2.00$ & $3.779\times10^{-4}$ & $4.00$\\
$64$ & $1.543\times10^{-2}$ & $2.00$ & $9.448\times10^{-5}$ & $4.00$\\
\bottomrule
\end{tabular}
\end{center}
Doubling $N$ halves the first-order error and quarters the second-order one,
exactly as the leading terms of Eq.~(\ref{eq:6-bch}) predict.  The ratios
approach their asymptotic values from above, which is the usual signature of
higher-order terms still contributing at small $N$.

\begin{notebox}
\textbf{Why this belongs in a chapter on mean fields.}  The immediate reason
is that Thouless' theorem is a statement about an exponential of one-body
operators, and the reason it is so simple is that those operators commute.
The longer-range reason is that the same expansion controls three later
developments.  In coupled-cluster theory the similarity-transformed
Hamiltonian $e^{-\hat{T}}\hat{H}e^{\hat{T}}$ is expanded by the
Baker-Hausdorff lemma
$e^{-X}Ye^{X}=Y-[X,Y]+\tfrac1{2!}[X,[X,Y]]-\cdots$, and terminates exactly at
the fourth commutator because $\hat{H}$ contains at most two-body terms.  In
time-dependent Hartree-Fock and in real-time methods generally, the propagator
is built by splitting.  And in quantum computing the Trotter formula is how a
Hamiltonian is turned into a circuit: each exponential
$e^{-\imath H_j\Delta}$ becomes a gate, Eq.~(\ref{eq:6-trotter}) fixes how
many of them are needed, and the trade-off between circuit depth and
accuracy is precisely the trade-off between $N$ and the error above.  A
first-order scheme needs $N=\bigO(t^2/\varepsilon)$ steps for an accuracy
$\varepsilon$; a second-order scheme needs only
$N=\bigO(t^{3/2}/\sqrt{\varepsilon})$.

The pieces $H_j$ that are split apart in that last application are the Pauli
strings of the Jordan-Wigner transformation of
Section~\ref{sec:jordanwigner}.  Writing $\hat{H}=\sum_{\alpha}h_{\alpha}
P_{\alpha}$ as in Eq.~(\ref{eq:3-paulisum}), each factor
$e^{-\imath h_{\alpha}P_{\alpha}\Delta}$ is a rotation generated by a single
Pauli string and has a standard circuit whose depth is set by the weight of
$P_{\alpha}$.  The number of strings is therefore the number of factors per
Trotter step, and their weights are the entangling cost of each -- which is
why Table~\ref{tab:3-paulicount} and Table~\ref{tab:4-paulicounts} were worth
computing.  Strings that commute with one another can be exponentiated
together with no splitting error at all, and grouping them is one of the main
practical optimisations.
\end{notebox}

%----------------------------------------------------------------------
\section{Stability of the Hartree-Fock solution}
\label{sec:hfstability}

\index{Hartree-Fock theory!stability}
The variational condition~(\ref{eq:6-hfcondition}) guarantees only that
$\bracket{c}{\hat{H}}{c}$ is \emph{stationary}.  A stationary point may be a
minimum, a maximum or a saddle, and which one it is matters: a Hartree-Fock
solution sitting on a saddle point is a poor reference for perturbation theory
or coupled cluster, because the corrections it generates are organised about
the wrong state.  Thouless' theorem lets us settle the question by going to
second order.

We ask whether
\begin{equation}
  \frac{\bracket{c'}{\hat{H}}{c'}}{\braket{c'}{c'}}
  \ \ge\ \bracket{c}{\hat{H}}{c} = E_0
  \label{eq:6-stabilitycondition}
\end{equation}
for every nearby determinant $\bket{c'}=\bket{c}+\bket{\delta c}$.  By
Thouless' theorem,
\begin{equation}
  \bket{c'} = \exp\left\{\sum_{a>F}\sum_{i\le F}\delta C_{ai}
    \hat{a}^{\dagger}_{a}\hat{a}_{i}\right\}\bket{c}
  = \left\{1+\sum_{ai}\delta C_{ai}\hat{a}^{\dagger}_{a}\hat{a}_{i}
    +\frac{1}{2!}\sum_{abij}\delta C_{ai}\delta C_{bj}
      \hat{a}^{\dagger}_{a}\hat{a}_{i}\hat{a}^{\dagger}_{b}\hat{a}_{j}
    +\cdots\right\}\bket{c},
  \label{eq:6-cprimeexpansion}
\end{equation}
with the amplitudes $\delta C$ small.  With intermediate normalisation
$\braket{c'}{c}=1$ the norm is
\begin{equation}
  \braket{c'}{c'} = 1 + \sum_{a>F}\sum_{i\le F}\vert\delta C_{ai}\vert^2
    + \bigO(\delta C^3).
  \label{eq:6-cprimenorm}
\end{equation}

\paragraph{The energy to second order.}
Expanding the numerator and using the Hartree-Fock
condition~(\ref{eq:6-hfcondition}) to kill the terms linear in $\delta C$,
\begin{align}
  \bracket{c'}{\hat{H}}{c'}
  &= \bracket{c}{\hat{H}}{c}
   + \sum_{ab>F}\sum_{ij\le F}\delta C^{*}_{ai}\delta C_{bj}
     \bracket{c}{\hat{a}^{\dagger}_{i}\hat{a}_{a}\hat{H}
       \hat{a}^{\dagger}_{b}\hat{a}_{j}}{c}
   \nonumber\\
  &\quad
   + \frac{1}{2!}\sum_{ab>F}\sum_{ij\le F}\delta C_{ai}\delta C_{bj}
     \bracket{c}{\hat{H}\hat{a}^{\dagger}_{a}\hat{a}_{i}
       \hat{a}^{\dagger}_{b}\hat{a}_{j}}{c}
   \nonumber\\
  &\quad
   + \frac{1}{2!}\sum_{ab>F}\sum_{ij\le F}\delta C^{*}_{ai}\delta C^{*}_{bj}
     \bracket{c}{\hat{a}^{\dagger}_{j}\hat{a}_{b}\hat{a}^{\dagger}_{i}
       \hat{a}_{a}\hat{H}}{c}
   + \cdots
  \label{eq:6-energysecondorder}
\end{align}
Each of these expectation values is evaluated with Wick's theorem,
Section~\ref{sec:wicktheorem}.  The first of them is the $1p$-$1h$ matrix
element computed in Section~\ref{sec:wickapps},
\begin{align}
  \bracket{c}{\{\hat{a}^{\dagger}_{i}\hat{a}_{a}\}\hat{H}
    \{\hat{a}^{\dagger}_{b}\hat{a}_{j}\}}{c}
  &= \sum_{pq}\bracket{p}{\hat{h}_0}{q}
     \bracket{c}{\{\hat{a}^{\dagger}_{i}\hat{a}_{a}\}
       \{\hat{a}^{\dagger}_{p}\hat{a}_{q}\}
       \{\hat{a}^{\dagger}_{b}\hat{a}_{j}\}}{c}
   \nonumber\\
  &\quad + \frac{1}{4}\sum_{pqrs}\langle pq\vert\hat{v}\vert rs\rangle
     \bracket{c}{\{\hat{a}^{\dagger}_{i}\hat{a}_{a}\}
       \{\hat{a}^{\dagger}_{p}\hat{a}^{\dagger}_{q}\hat{a}_{s}\hat{a}_{r}\}
       \{\hat{a}^{\dagger}_{b}\hat{a}_{j}\}}{c},
  \label{eq:6-onephonemat}
\end{align}
and collecting the surviving contractions gives
\begin{equation}
  E_0\sum_{ai}\vert\delta C_{ai}\vert^2
  + \sum_{ai}\vert\delta C_{ai}\vert^2
    \left(\varepsilon_a-\varepsilon_i\right)
  - \sum_{ijab}\langle aj\vert\hat{v}\vert bi\rangle
    \,\delta C^{*}_{ai}\delta C_{bj}.
  \label{eq:6-firstblock}
\end{equation}
The third line of Eq.~(\ref{eq:6-energysecondorder}) is the Hermitian
conjugate of the second, since
$\bracket{a}{\hat{A}}{b}=\left(\bracket{b}{\hat{A}^{\dagger}}{a}\right)^{*}$
and $\hat{H}$ is Hermitian, and the second evaluates to
\begin{equation}
  \frac{1}{2}\sum_{ijab}\langle ab\vert\hat{v}\vert ij\rangle
    \,\delta C^{*}_{ai}\delta C^{*}_{bj} .
  \label{eq:6-secondblock}
\end{equation}

\paragraph{The stability matrix.}
Define the two antisymmetrised blocks
\begin{equation}
  A_{ai,bj} = -\langle aj\vert\hat{v}\vert bi\rangle_{\mathrm{AS}},
  \qquad
  B_{ai,bj} = \langle ab\vert\hat{v}\vert ij\rangle_{\mathrm{AS}},
  \label{eq:6-ABblocks}
\end{equation}
together with the diagonal matrix of unperturbed excitation energies
\begin{equation}
  \Delta_{ai,bj}
    = \left(\varepsilon_a-\varepsilon_i\right)
      \delta_{ab}\delta_{ij}.
  \label{eq:6-Delta}
\end{equation}
Then Eqs.~(\ref{eq:6-firstblock}) and~(\ref{eq:6-secondblock}) combine into
\begin{align}
  \bracket{c'}{\hat{H}}{c'}
  &= \left(1+\sum_{ai}\vert\delta C_{ai}\vert^2\right)
     \bracket{c}{\hat{H}}{c}
   + \sum_{ai}\vert\delta C_{ai}\vert^2
     \left(\varepsilon^{\mathrm{HF}}_a-\varepsilon^{\mathrm{HF}}_i\right)
   + \sum_{ijab}A_{ai,bj}\,\delta C^{*}_{ai}\delta C_{bj}
   \nonumber\\
  &\quad
   + \frac{1}{2}\sum_{ijab}B^{*}_{ai,bj}\,\delta C_{ai}\delta C_{bj}
   + \frac{1}{2}\sum_{ijab}B_{ai,bj}\,
       \delta C^{*}_{ai}\delta C^{*}_{bj}
   + \bigO(\delta C^3),
  \label{eq:6-energyquadratic}
\end{align}
so that, dividing by the norm~(\ref{eq:6-cprimenorm}),
\begin{equation}
  \frac{\bracket{c'}{\hat{H}}{c'}}{\braket{c'}{c'}}
  = E_0 + \frac{\Delta E}
    {1+\sum_{ai}\vert\delta C_{ai}\vert^2},
  \qquad
  \Delta E = \frac{1}{2}\bracket{\chi}{\hat{M}}{\chi},
  \label{eq:6-deltaE}
\end{equation}
with the composite vector and matrix
\begin{equation}
  \chi = \begin{bmatrix}\delta C & \delta C^{*}\end{bmatrix}^{T},
  \qquad
  \hat{M} = \begin{pmatrix}
    \Delta + A & B\\
    B^{*} & \Delta + A^{*}
  \end{pmatrix}.
  \label{eq:6-stabilitymatrix}
\end{equation}

The conclusion is now immediate.  The denominator in
Eq.~(\ref{eq:6-deltaE}) is positive, so the Hartree-Fock solution is a local
minimum if and only if
\begin{equation}
  \boxed{\;
  \Delta E = \frac{1}{2}\bracket{\chi}{\hat{M}}{\chi} \ \ge\ 0
  \quad\text{for every }\chi
  \;}
  \label{eq:6-stabilitycriterion}
\end{equation}
that is, if and only if every eigenvalue of $\hat{M}$ is non-negative.  A
necessary but not sufficient condition, obtained by restricting $\chi$ to the
upper block, is that
\begin{equation}
  \left(\varepsilon_a-\varepsilon_i\right)\delta_{ab}\delta_{ij}
  + A_{ai,bj} \ \ge\ 0
  \label{eq:6-necessary}
\end{equation}
as a matrix.  This is cheap to test and is the usual first check.

\begin{notebox}
\textbf{The stability matrix is the RPA matrix.}  The upper-left block
$\Delta+A$ is exactly the matrix diagonalised in the Tamm-Dancoff
approximation: it is the Hamiltonian restricted to the $1p$-$1h$ space,
measured from the reference energy.  The full $\hat{M}$, with its off-diagonal
$B$ blocks, is the matrix of the random-phase approximation.  So the statement
``the Hartree-Fock solution is stable'' and the statement ``the RPA
frequencies are real'' are one and the same, and an imaginary RPA frequency is
not a numerical accident but a signal that the mean field itself is wrong.
This is why Chapter~\ref{chap:fci} could describe TDA and RPA as truncations
of full configuration interaction and why they belong immediately after this
chapter: they are the second derivative of the Hartree-Fock energy surface.
\end{notebox}

%----------------------------------------------------------------------
\section{An explicit instability: the Lipkin model}
\label{sec:lipkininstability}

The Lipkin model of Section~\ref{sec:lipkin} is the standard illustration,
because everything can be computed in closed form.  Take $W=0$, so that
\begin{equation}
  \hat{H} = \varepsilon\hat{J}_z
    + \frac{V}{2}\left(\hat{J}_+^2+\hat{J}_-^2\right),
  \label{eq:6-lipkinH}
\end{equation}
with $N$ particles in two $N$-fold degenerate levels separated by
$\varepsilon$.  The natural mean-field ansatz rotates every single-particle
state by the same angle $\alpha$ in the two-level space,
\begin{equation}
  \hat{b}^{\dagger}_{p} = \cos\!\left(\tfrac{\alpha}{2}\right)
    \hat{a}^{\dagger}_{p-}
    + \sin\!\left(\tfrac{\alpha}{2}\right)\hat{a}^{\dagger}_{p+},
  \label{eq:6-lipkinrotation}
\end{equation}
which is a Thouless rotation with $C_{ai}=\tan(\alpha/2)\,\delta_{pq}$.  The
resulting energy is the classic result
\begin{equation}
  \frac{E(\alpha)}{N} = -\frac{\varepsilon}{2}\cos\alpha
    - \frac{\varepsilon\chi}{4}\sin^2\alpha,
  \qquad
  \chi = \frac{\vert V\vert (N-1)}{\varepsilon},
  \label{eq:6-lipkinmeanfield}
\end{equation}
where $\chi$ is the dimensionless coupling.  Differentiating,
\begin{equation}
  \frac{1}{N}\frac{\ud{E}}{\ud{\alpha}}
   = \frac{\varepsilon}{2}\sin\alpha
     \left(1-\chi\cos\alpha\right),
  \label{eq:6-lipkinderivative}
\end{equation}
so there are two stationary points: the \emph{spherical} solution
$\alpha=0$, which exists always, and a \emph{deformed} solution
$\cos\alpha=1/\chi$, which exists only for $\chi\ge1$.  The second
derivative at $\alpha=0$ is
\begin{equation}
  \frac{1}{N}\left.\frac{\ud{^2E}}{\ud{\alpha^2}}\right\vert_{\alpha=0}
   = \frac{\varepsilon}{2}\left(1-\chi\right),
  \label{eq:6-lipkinsecond}
\end{equation}
which changes sign at $\chi=1$.  The spherical solution turns from a minimum
into a maximum exactly where the deformed solution appears, which is the
signature of a continuous symmetry-breaking transition.

The stability matrix says the same thing without any reference to the
parametrisation~(\ref{eq:6-lipkinrotation}).  Because the interaction
in Eq.~(\ref{eq:6-lipkinH}) changes $J_z$ by two units, it has no matrix
element within the $1p$-$1h$ space, so $A=0$ and $\Delta+A=\varepsilon\bm{I}$.
The block $B$ has matrix elements $\pm V$ connecting different pairs, and its
largest eigenvalue is $\vert V\vert(N-1)$.  The lowest eigenvalue of
$\hat{M}$ is therefore
\begin{equation}
  \lambda_{\min} = \varepsilon - \vert V\vert(N-1)
    = \varepsilon\left(1-\chi\right),
  \label{eq:6-lipkinlowest}
\end{equation}
which is exactly Eq.~(\ref{eq:6-lipkinsecond}) up to the factor $N/2$.
The program \texttt{hartreefock.py} builds $\hat{M}$ numerically from the
Hamiltonian in the determinant basis, for $N=4$ and $\varepsilon=1$:

\begin{table}[h]
\centering
\renewcommand{\arraystretch}{1.25}
\setlength{\tabcolsep}{9pt}
\begin{tabular}{rrrrrl}
\toprule
$\chi$ & $\lambda_{\min}(\hat{M})$ & $\alpha_{\min}$
  & $E_{\mathrm{mf}}$ & $E_{\mathrm{exact}}$ & solution\\
\midrule
$0.25$ & $\phantom{-}0.75$ & $0.000000$ & $-2.000000$ & $-2.020726$ & spherical\\
$0.50$ & $\phantom{-}0.50$ & $0.000000$ & $-2.000000$ & $-2.081666$ & spherical\\
$0.90$ & $\phantom{-}0.10$ & $0.000000$ & $-2.000000$ & $-2.253886$ & spherical\\
$1.00$ & $\phantom{-}0.00$ & $0.000000$ & $-2.000000$ & $-2.309401$ & critical\\
$1.50$ & $-0.50$ & $0.841069$ & $-2.166667$ & $-2.645751$ & deformed\\
$2.00$ & $-1.00$ & $1.047198$ & $-2.500000$ & $-3.055050$ & deformed\\
\bottomrule
\end{tabular}
\caption{Stability of the spherical Hartree-Fock solution of the Lipkin model
with $N=4$, $\varepsilon=1$ and $W=0$.  The lowest eigenvalue of the
stability matrix crosses zero exactly at $\chi=1$, where the deformed
minimum appears.  $E_{\mathrm{mf}}$ is the mean-field energy at the minimum
and $E_{\mathrm{exact}}$ the exact ground-state energy.}
\label{tab:lipkinstability}
\end{table}

Two lessons are worth drawing.  The first is that a solution can satisfy the
Hartree-Fock equations perfectly and still be useless: for $\chi>1$ the
spherical state is a genuine stationary point, obeys Brillouin's theorem, and
is a saddle.  The second is that the deformed solution buys a lower energy by
breaking a symmetry the exact ground state does not break -- at $\chi=2$ it
gains $0.5$ but is still $0.55$ above the exact answer.  Symmetry breaking in
a mean field is a device for capturing correlations cheaply, and the symmetry
must eventually be restored, by projection or by a method that never broke it.

%----------------------------------------------------------------------
\section{When the mean field does nothing}
\label{sec:hfnothing}

It is instructive to see the opposite extreme.  Apply the same machinery to
the pairing model of Section~\ref{sec:pairingmodel}, with four doubly
degenerate levels and four particles.  The program reports

\begin{center}
\renewcommand{\arraystretch}{1.2}
\begin{tabular}{rllll}
\toprule
$g$ & $E$(reference) & $E$(Hartree-Fock) & $E$(exact) & iterations\\
\midrule
$0.25$ & $1.75000000$ & $1.75000000$ & $1.73045985$ & $2$\\
$0.50$ & $1.50000000$ & $1.50000000$ & $1.41677428$ & $2$\\
$1.00$ & $1.00000000$ & $1.00000000$ & $0.63554847$ & $2$\\
\bottomrule
\end{tabular}
\end{center}

Hartree-Fock gains nothing at all.  The reason is the seniority conservation
of Section~\ref{sec:fcipairing}: the pairing interaction moves complete pairs
and cannot connect the reference to a $1p$-$1h$ state, so
$\bracket{\Phi_0}{\hat{H}}{\Phi_i^a}=0$ identically, the Fock matrix is
diagonal from the first iteration, and the self-consistency loop terminates
without having moved.  The entire correlation energy -- $36\%$ of the
reference energy at $g=1$ -- must come from the doubles, which is precisely
what Table~\ref{tab:pairing} and Section~\ref{sec:citruncations} showed.

The moral is that a mean field is only as useful as the one-body physics the
interaction actually contains.  Where an interaction can be well approximated
by an average field -- atoms, molecules, closed-shell nuclei -- Hartree-Fock
does most of the work.  Where the interaction is intrinsically a
pair-correlation effect -- pairing, superconductivity, the strongly coupled
Hubbard model of Section~\ref{sec:fcihubbard} -- it does none, and one must
either generalise the mean field, as Bardeen, Cooper and Schrieffer did by
allowing the number of particles to fluctuate, or go beyond it.

%----------------------------------------------------------------------
\section{The infinite homogeneous electron gas}
\label{sec:electrongas}

\index{electron gas}
The electron gas is very likely the only realistic model of an interacting
many-particle system for which the Hartree-Fock equations can be solved in
closed form.  Better still, the total energy and several other properties
follow analytically to first order in the interaction, so the whole chapter
can be tested against an exact mean-field answer.

The model is defined by three assumptions:
\begin{itemize}
  \item the electrons feel no external force except the attraction of a
    uniform, stationary background of positive charge;
  \item the system as a whole is neutral;
  \item there are $N$ electrons in a cubic box of side $L$ and volume
    $\Omega=L^3$, containing a uniform positive charge of density $Ne/\Omega$.
\end{itemize}
It is a very good approximation for the valence electrons of a metal, and the
three-dimensional gas is the cornerstone of the local-density approximation of
density-functional theory.  Two-dimensional electron fluids occur on metal and
liquid-helium surfaces and at metal-oxide-semiconductor interfaces.  For all
its simplicity, the correlations of this system have resisted a satisfactory
first-principles description by anything except quantum Monte Carlo; the
diffusion Monte Carlo calculations of Ceperley, and of Ceperley and Tanatar,
remain the benchmark in three and in two dimensions.  At low density the
electrons localise into a lattice -- Wigner crystallisation, a direct
consequence of the long range of the repulsion -- while at high density the
system behaves as a liquid.  The infinite range of the Coulomb interaction
also generates finite-size effects that are absent in nuclear or neutron
matter, and that any finite calculation must handle with care.

\subsection{Plane waves and periodic boundary conditions}
\label{sec:egbasis}

The system is homogeneous, so the single-particle states are plane waves
normalised in the box,
\begin{equation}
  \psi_{\bm{k}\sigma}(\bm{r})
   = \frac{1}{\sqrt{\Omega}}e^{\imath\bm{k}\cdot\bm{r}}\,\xi_{\sigma},
  \qquad
  \xi_{+1/2}=\begin{pmatrix}1\\0\end{pmatrix},
  \quad
  \xi_{-1/2}=\begin{pmatrix}0\\1\end{pmatrix},
  \label{eq:6-planewave}
\end{equation}
and periodic boundary conditions restrict the wave numbers to
\begin{equation}
  k_i = \frac{2\pi n_i}{L},
  \qquad i=x,y,z,
  \qquad n_i = 0,\pm1,\pm2,\dots
  \label{eq:6-pbc}
\end{equation}
The limit $L\to\infty$ is taken after the expectation values have been
computed.  The kinetic energy operator is diagonal,
\begin{equation}
  \hat{T} = \sum_{\bm{k}\sigma}\frac{\hbar^2k^2}{2m}
    \hat{a}^{\dagger}_{\bm{k}\sigma}\hat{a}_{\bm{k}\sigma},
  \label{eq:6-kinetic}
\end{equation}
which is the same structure as the momentum-space Hubbard model of
Section~\ref{sec:fcihubbard}: the plane-wave determinant is automatically the
Hartree-Fock solution of a translationally invariant problem, since momentum
conservation forces $f_{ai}=0$.  Everything in this section is therefore a
calculation \emph{in} the Hartree-Fock basis, not a search for it.

The full Hamiltonian has three pieces,
\begin{equation}
  \hat{H} = \hat{H}_{\mathrm{el}} + \hat{H}_{b} + \hat{H}_{\mathrm{el}-b},
  \label{eq:6-threepieces}
\end{equation}
namely the electrons,
\begin{equation}
  \hat{H}_{\mathrm{el}} = \sum_{i=1}^{N}\frac{p_i^2}{2m}
    + \frac{e^2}{2}\sum_{i\ne j}
      \frac{e^{-\mu\vert\bm{r}_i-\bm{r}_j\vert}}
           {\vert\bm{r}_i-\bm{r}_j\vert},
  \label{eq:6-Hel}
\end{equation}
the background with itself,
\begin{equation}
  \hat{H}_{b} = \frac{e^2}{2}\iint \ud{\bm{r}}\ud{\bm{r}'}
    \frac{n(\bm{r})n(\bm{r}')e^{-\mu\vert\bm{r}-\bm{r}'\vert}}
         {\vert\bm{r}-\bm{r}'\vert},
  \label{eq:6-Hb}
\end{equation}
with $n(\bm{r})=N/\Omega$, and the electrons with the background,
\begin{equation}
  \hat{H}_{\mathrm{el}-b} = -\frac{e^2}{2}\sum_{i=1}^{N}\int \ud{\bm{r}}
    \frac{n(\bm{r})e^{-\mu\vert\bm{r}-\bm{x}_i\vert}}
         {\vert\bm{r}-\bm{x}_i\vert}.
  \label{eq:6-Helb}
\end{equation}
The factor $e^{-\mu r}$ is a convergence device: each integral is computed at
finite $\mu$, and the limit $\mu\to0$ is taken at the end.  Individually the
three terms diverge as $\mu^{-2}$,
\begin{equation}
  \hat{H}_{b} = \frac{e^2}{2}\frac{N^2}{\Omega}\frac{4\pi}{\mu^2},
  \qquad
  \hat{H}_{\mathrm{el}-b} = -e^2\frac{N^2}{\Omega}\frac{4\pi}{\mu^2},
  \label{eq:6-divergences}
\end{equation}
and the direct term of the electron-electron repulsion supplies exactly the
missing $+\tfrac12$.  The three divergences cancel, and what survives is the
kinetic energy and the exchange term.  This cancellation is special to the
Coulomb interaction and to a neutral system; for a short-ranged interaction,
nuclear forces for example, both direct and exchange terms must be kept.  What
remains is
\begin{equation}
  \hat{H} = \hat{H}_0 + \hat{H}_I,
  \qquad
  \hat{H}_0 = \sum_{\bm{k}\sigma}\frac{\hbar^2k^2}{2m}
    \hat{a}^{\dagger}_{\bm{k}\sigma}\hat{a}_{\bm{k}\sigma},
  \label{eq:6-egH0}
\end{equation}
\begin{equation}
  \hat{H}_I = \frac{e^2}{2\Omega}\sum_{\sigma_1\sigma_2}
    \sum_{\bm{q}\ne0,\bm{k},\bm{p}}\frac{4\pi}{q^2}\,
    \hat{a}^{\dagger}_{\bm{k}+\bm{q},\sigma_1}
    \hat{a}^{\dagger}_{\bm{p}-\bm{q},\sigma_2}
    \hat{a}_{\bm{p}\sigma_2}\hat{a}_{\bm{k}\sigma_1},
  \label{eq:6-egHI}
\end{equation}
with the $\bm{q}=0$ term -- the one that diverged -- removed by the
background.

\subsection{The interaction in momentum space}
\label{sec:egmomentum}

Equation~(\ref{eq:6-egHI}) deserves a derivation, since the same steps apply
to any interaction depending only on the relative distance.  Starting from
\begin{equation}
  \hat{V} = \frac{1}{2}\sum_{\substack{\sigma_p\sigma_q\\ \sigma_r\sigma_s}}
    \sum_{\substack{\bm{k}_p\bm{k}_q\\ \bm{k}_r\bm{k}_s}}
    \langle \bm{k}_p\sigma_p,\bm{k}_q\sigma_q\vert\hat{v}\vert
      \bm{k}_r\sigma_r,\bm{k}_s\sigma_s\rangle\,
    \hat{a}^{\dagger}_{\bm{k}_p\sigma_p}
    \hat{a}^{\dagger}_{\bm{k}_q\sigma_q}
    \hat{a}_{\bm{k}_s\sigma_s}\hat{a}_{\bm{k}_r\sigma_r},
  \label{eq:6-Vgeneral}
\end{equation}
insert the plane waves~(\ref{eq:6-planewave}).  The spin functions are
orthonormal and $\hat{v}$ is spin independent, so
\begin{equation}
  \langle \bm{k}_p\sigma_p,\bm{k}_q\sigma_q\vert\hat{v}\vert
    \bm{k}_r\sigma_r,\bm{k}_s\sigma_s\rangle
  = \frac{\delta_{\sigma_p\sigma_r}\delta_{\sigma_q\sigma_s}}{\Omega^2}
    \iint e^{-\imath\bm{k}_p\cdot\bm{r}_p}
          e^{-\imath\bm{k}_q\cdot\bm{r}_q}\,v(r)\,
          e^{\imath\bm{k}_r\cdot\bm{r}_p}
          e^{\imath\bm{k}_s\cdot\bm{r}_q}
      \ud{\bm{r}_p}\ud{\bm{r}_q}.
  \label{eq:6-planewaveme}
\end{equation}
Change variables to $\bm{r}=\bm{r}_p-\bm{r}_q$ at fixed $\bm{r}_q$; the limits
are unchanged, and the exponentials separate,
\begin{equation}
  = \frac{\delta_{\sigma_p\sigma_r}\delta_{\sigma_q\sigma_s}}{\Omega^2}
    \int v(r)e^{\imath(\bm{k}_r-\bm{k}_p)\cdot\bm{r}}\ud{\bm{r}}
    \int e^{\imath(\bm{k}_s-\bm{k}_q+\bm{k}_r-\bm{k}_p)\cdot\bm{r}_q}
      \ud{\bm{r}_q}.
  \label{eq:6-separated}
\end{equation}
The second integral is a Kronecker delta enforcing momentum conservation,
\begin{equation}
  \langle \bm{k}_p\sigma_p,\bm{k}_q\sigma_q\vert\hat{v}\vert
    \bm{k}_r\sigma_r,\bm{k}_s\sigma_s\rangle
  = \frac{\delta_{\sigma_p\sigma_r}\delta_{\sigma_q\sigma_s}}{\Omega}
    \,\delta_{\bm{k}_p+\bm{k}_q,\,\bm{k}_r+\bm{k}_s}
    \int v(r)e^{\imath(\bm{k}_r-\bm{k}_p)\cdot\bm{r}}\ud{\bm{r}} .
  \label{eq:6-momentumconservation}
\end{equation}
The matrix element vanishes unless
$\bm{k}_p+\bm{k}_q=\bm{k}_r+\bm{k}_s$.  Using this to eliminate one summation
variable and writing $\bm{p}=\bm{k}_p+\bm{k}_q-\bm{k}_r$,
$\bm{k}=\bm{k}_r$ and $\bm{q}=\bm{k}_p-\bm{k}_r$,
\begin{equation}
  \hat{V} = \frac{1}{2\Omega}\sum_{\sigma\sigma'}
    \sum_{\bm{k}\bm{p}\bm{q}}
    \left[\int v(r)e^{-\imath\bm{q}\cdot\bm{r}}\ud{\bm{r}}\right]
    \hat{a}^{\dagger}_{\bm{k}+\bm{q},\sigma}
    \hat{a}^{\dagger}_{\bm{p}-\bm{q},\sigma'}
    \hat{a}_{\bm{p}\sigma'}\hat{a}_{\bm{k}\sigma}.
  \label{eq:6-Vmomentum}
\end{equation}
For the screened Coulomb interaction the bracket is
$4\pi e^2/(\mu^2+q^2)$, which becomes $4\pi e^2/q^2$ as $\mu\to0$, giving
Eq.~(\ref{eq:6-egHI}).  The structure -- a momentum transfer $\bm{q}$ shared
between two particles -- is exactly the Hubbard interaction of
Eq.~(\ref{eq:5-hubbardmom}), with $4\pi e^2/q^2$ in place of the constant $U$.

\subsection{The reference energy}
\label{sec:egreference}

With $\bket{\Phi_0}$ the filled Fermi sea, the kinetic energy follows from the
contractions of Section~\ref{sec:wickapps},
\begin{equation}
  \bracket{\Phi_0}{\hat{T}}{\Phi_0}
   = \sum_{i\le F}\frac{\hbar^2k_i^2}{2m}\cdot 2
   = \frac{2\Omega}{(2\pi)^3}\frac{\hbar^2}{2m}
     \int_0^{k_F}4\pi k^4\ud{k}
   = \frac{\hbar^2\Omega}{10\pi^2 m}k_F^5,
  \label{eq:6-egkinetic}
\end{equation}
the factor two being the spin degeneracy.  The same density of states fixes
the particle number,
\begin{equation}
  N = \frac{2\Omega}{(2\pi)^3}\cdot\frac{4}{3}\pi k_F^3
    = \frac{\Omega}{3\pi^2}k_F^3
  \qquad\Longrightarrow\qquad
  k_F = \left(\frac{3\pi^2N}{\Omega}\right)^{1/3},
  \label{eq:6-kF}
\end{equation}
so that
\begin{equation}
  \frac{\bracket{\Phi_0}{\hat{T}}{\Phi_0}}{N}
   = \frac{3}{5}\frac{\hbar^2k_F^2}{2m}
   = \frac{\hbar^2(3\pi^2)^{5/3}}{10\pi^2 m}\rho^{2/3},
  \qquad \rho = \frac{N}{\Omega}.
  \label{eq:6-egkineticdensity}
\end{equation}
For the interaction, only two classes of contraction survive: either
$\bm{k}+\bm{q}=\bm{p}$ with $\sigma=\sigma'$, or $\bm{q}=0$.  The first is the
exchange term, the second the direct one, and
\begin{equation}
  \bracket{\Phi_0}{\hat{V}}{\Phi_0}
   = \frac{1}{2\Omega}\left(N^2\int v(r)\ud{\bm{r}}
     - N\sum_{\bm{q}}\int v(r)e^{-\imath\bm{q}\cdot\bm{r}}\ud{\bm{r}}\right),
  \label{eq:6-egpotential}
\end{equation}
valid for any interaction depending only on $r$.  For the electron gas the
first term is cancelled by the background, as we saw, and only the exchange
term remains.

\subsection{The single-particle energy}
\label{sec:egsingleparticle}

The Hartree-Fock single-particle energy is the kinetic energy plus the
exchange self-energy,
\begin{equation}
  \varepsilon_{k}^{\mathrm{HF}}
   = \frac{\hbar^2k^2}{2m}
   - \frac{e^2}{\Omega^2}\sum_{p\le k_F}
     \int \ud{\bm{r}}\,e^{\imath(\bm{p}-\bm{k})\cdot\bm{r}}
     \int \ud{\bm{r}'}
     \frac{e^{\imath(\bm{k}-\bm{p})\cdot\bm{r}'}}
          {\vert\bm{r}-\bm{r}'\vert},
  \label{eq:6-egsp}
\end{equation}
and we shall show that it evaluates to
\begin{equation}
  \boxed{\;
  \varepsilon_{k}^{\mathrm{HF}}
   = \frac{\hbar^2k^2}{2m}
   - \frac{e^2k_F}{2\pi}\left[2
     + \frac{k_F^2-k^2}{kk_F}
       \ln\left\vert\frac{k+k_F}{k-k_F}\right\vert\right].
  \;}
  \label{eq:6-egspclosed}
\end{equation}

Insert the convergence factor $e^{-\mu\vert\bm{r}-\bm{r}'\vert}$ and replace
the sum by an integral, $\sum_{\bm{p}}\to\Omega(2\pi)^{-3}\int\ud{\bm{p}}$,
\begin{equation}
  \lim_{\mu\to0}\frac{e^2}{\Omega(2\pi)^3}
    \iint \ud{\bm{r}}\ud{\bm{r}'}
    \frac{e^{-\mu\vert\bm{r}-\bm{r}'\vert}}{\vert\bm{r}-\bm{r}'\vert}
    \int \ud{\bm{p}}\,
      e^{\imath(\bm{p}-\bm{k})\cdot(\bm{r}-\bm{r}')} .
  \label{eq:6-egstep1}
\end{equation}
Change variables to $\bm{x}=\bm{r}-\bm{r}'$ and $\bm{y}=\bm{r}'$.  The
integrand no longer depends on $\bm{y}$, whose integral gives $\Omega$, and
the $\bm{x}$ integral in spherical coordinates is
\begin{equation}
  \int \ud{\bm{x}}\,e^{\imath(\bm{p}-\bm{k})\cdot\bm{x}}
    \frac{e^{-\mu x}}{x}
  = 4\pi\int_0^{\infty}
    \frac{\sin\left(\vert\bm{p}-\bm{k}\vert x\right)}
         {\vert\bm{p}-\bm{k}\vert}e^{-\mu x}\ud{x}
  = \frac{4\pi}{\mu^2+\vert\bm{p}-\bm{k}\vert^2}.
  \label{eq:6-egstep2}
\end{equation}
The limit $\mu\to0$ is now harmless, and we are left with
\begin{equation}
  \frac{e^2}{2\pi^2}\int_{p\le k_F}
    \frac{\ud{\bm{p}}}{\vert\bm{p}-\bm{k}\vert^2}
  = \frac{e^2}{\pi}\int_0^{k_F}p^2\ud{p}
    \int_{-1}^{1}\frac{\ud{u}}{p^2+k^2-2pku},
  \label{eq:6-egstep3}
\end{equation}
having set $u=\cos\theta$.  The angular integral is elementary,
\begin{equation}
  \int_{-1}^{1}\frac{\ud{u}}{p^2+k^2-2pku}
   = \frac{1}{pk}\ln\left\vert\frac{p+k}{p-k}\right\vert,
  \label{eq:6-egangular}
\end{equation}
so the exchange term becomes
\begin{equation}
  \frac{e^2}{k\pi}\int_0^{k_F}p\,
    \ln\left\vert\frac{p+k}{p-k}\right\vert\ud{p}
  = \frac{e^2}{k\pi}\left[kk_F
    + \frac{k_F^2-k^2}{2}
      \ln\left\vert\frac{k+k_F}{k-k_F}\right\vert\right],
  \label{eq:6-egstep4}
\end{equation}
where the remaining integral is done by substituting $x=p+k$ and $y=p-k$.
Rearranging gives Eq.~(\ref{eq:6-egspclosed}).

\subsection{Dimensionless form, band width and effective mass}
\label{sec:egdimensionless}

Introduce the radius $r_0$ of the sphere whose volume is the volume per
electron,
\begin{equation}
  \rho = \frac{k_F^3}{3\pi^2} = \frac{3}{4\pi r_0^3}
  \qquad\Longrightarrow\qquad
  k_F = \frac{(9\pi/4)^{1/3}}{r_0} = \frac{1.9192}{r_0},
  \label{eq:6-rs}
\end{equation}
the Bohr radius $a_0=\hbar^2/me^2$ and the dimensionless density parameter
$r_s=r_0/a_0$, which lies between about $2$ and $6$ for most metals.  With
$x=k/k_F$ and
\begin{equation}
  F(x) = \frac{1}{2} + \frac{1-x^2}{4x}
    \ln\left\vert\frac{1+x}{1-x}\right\vert,
  \label{eq:6-Fx}
\end{equation}
Equation~(\ref{eq:6-egspclosed}) becomes
\begin{equation}
  \varepsilon_{k}^{\mathrm{HF}}
   = \frac{\hbar^2k^2}{2m} - \frac{2e^2k_F}{\pi}F(x),
  \qquad
  \frac{\varepsilon_{k}^{\mathrm{HF}}}{\varepsilon_0^{F}}
   = x^2 - \frac{4}{\pi k_F a_0}F(x)
   = x^2 - 0.663\,r_s\,F(x),
  \label{eq:6-egdimensionless}
\end{equation}
where $\varepsilon_0^{F}=\hbar^2k_F^2/2m$ is the free energy at the Fermi
surface.  The function $F$ is smooth except at $x=1$, where it takes the value
$\tfrac12$ but has an infinite derivative; at $x=0$ it equals $1$.  The
program \texttt{hartreefock.py} tabulates the result for $r_s=4$:

\begin{table}[h]
\centering
\renewcommand{\arraystretch}{1.2}
\setlength{\tabcolsep}{12pt}
\begin{tabular}{rrrr}
\toprule
$x=k/k_F$ & $F(x)$ & $\varepsilon_k^{\mathrm{HF}}/\varepsilon_0^{F}$
  & $\varepsilon_k^{(0)}/\varepsilon_0^{F}$\\
\midrule
$0.00$ & $1.000000$ & $-2.652000$ & $0.0000$\\
$0.25$ & $0.978899$ & $-2.533540$ & $0.0625$\\
$0.50$ & $0.911980$ & $-2.168570$ & $0.2500$\\
$0.75$ & $0.783779$ & $-1.516081$ & $0.5625$\\
$1.00$ & $0.500000$ & $-0.326000$ & $1.0000$\\
$1.25$ & $0.252812$ & $\phantom{-}0.892042$ & $1.5625$\\
$1.50$ & $0.164700$ & $\phantom{-}1.813214$ & $2.2500$\\
$2.00$ & $0.088020$ & $\phantom{-}3.766570$ & $4.0000$\\
\bottomrule
\end{tabular}
\caption{The Hartree-Fock single-particle energy of the three-dimensional
electron gas at $r_s=4$, compared with the free-particle energy.  The
exchange term lowers every state and lowers the bottom of the band far more
than the top.}
\label{tab:eghf}
\end{table}

\paragraph{Band width.}
The width of the occupied band is
\begin{equation}
  \Delta\varepsilon^{\mathrm{HF}}
   = \varepsilon_{k_F}^{\mathrm{HF}} - \varepsilon_{0}^{\mathrm{HF}} .
  \label{eq:6-bandwidth}
\end{equation}
At $r_s=4$ the two entries of Table~\ref{tab:eghf} give $-0.326$ and
$-2.652$, hence $\Delta\varepsilon^{\mathrm{HF}}=2.326$ in units of
$\varepsilon_0^{F}$.  In general, using $F(1)=\tfrac12$ and $F(0)=1$,
\begin{equation}
  \Delta\varepsilon^{\mathrm{HF}}
   = 1 + \frac{0.663}{2}r_s
   = 1 + 0.3315\,r_s ,
  \label{eq:6-bandwidthformula}
\end{equation}
which reproduces $2.326$ at $r_s=4$.  The free gas has a band width of exactly
one, so exchange more than doubles it -- an effect of the right sign but, as
experiment shows, considerably too large.

\paragraph{Effective mass and level density.}
Expanding about the Fermi surface,
\begin{equation}
  \varepsilon_k^{\mathrm{HF}} = \varepsilon_{k_F}^{\mathrm{HF}}
    + \left(\frac{\partial\varepsilon_k^{\mathrm{HF}}}
      {\partial k}\right)_{k_F}(k-k_F) + \cdots
  \label{eq:6-taylor}
\end{equation}
and comparing with the free-particle expansion
$\varepsilon_k^{(0)}=\hbar^2k_F^2/2m+(\hbar^2k_F/m)(k-k_F)+\cdots$ defines the
effective mass
\begin{equation}
  m^{*}_{\mathrm{HF}}
   \equiv \hbar^2k_F\left(
     \frac{\partial\varepsilon_k^{\mathrm{HF}}}{\partial k}
   \right)^{-1}_{k_F}
   = \frac{2m}{2-0.663\,r_s\,F'(1)} .
  \label{eq:6-effectivemass}
\end{equation}
Here is the difficulty: $F'(x)$ diverges logarithmically as $x\to1$, so the
slope of $\varepsilon^{\mathrm{HF}}_k$ is infinite at the Fermi surface and
$m^{*}_{\mathrm{HF}}\to0$.  Approaching from below, the program finds
$m^{*}/m = 0.256, 0.185, 0.144, 0.118, 0.100$ at
$1-x=10^{-2},\dots,10^{-6}$: the ratio creeps towards zero like
$1/\ln(1-x)$.  Since the level density is
\begin{equation}
  n(\varepsilon) = \frac{\Omega k^2}{2\pi^2}
    \left(\frac{\partial\varepsilon}{\partial k}\right)^{-1},
  \label{eq:6-leveldensity}
\end{equation}
it vanishes at the Fermi surface.

This is qualitatively wrong.  Metals have a finite density of states at the
Fermi level -- it is what gives them a linear specific heat and a finite Pauli
susceptibility.  The culprit is the $1/q^2$ singularity of the bare Coulomb
interaction at small momentum transfer, which Hartree-Fock treats without any
screening.  In a real metal the other electrons rearrange to screen it, and
the effective interaction falls off exponentially beyond the screening length.
Summing the diagrams that produce that screening is precisely what the
random-phase approximation does, and it restores a finite effective mass.  So
the electron gas illustrates both what a mean field does well and where it
must fail, and it points directly to the next method.

\subsection{The energy per electron}
\label{sec:egenergy}

Collecting the kinetic energy~(\ref{eq:6-egkineticdensity}) and the exchange
term, and measuring energies in Rydberg, $e^2/2a_0=13.606$~eV, the
Hartree-Fock energy per electron is
\begin{equation}
  \boxed{\;
  \frac{E_0}{N} = \frac{e^2}{2a_0}
    \left[\frac{2.21}{r_s^2}-\frac{0.916}{r_s}\right].
  \;}
  \label{eq:6-egenergy}
\end{equation}
The two coefficients are not empirical:
\begin{equation}
  \frac{3}{5}\left(\frac{9\pi}{4}\right)^{2/3} = 2.2099,
  \qquad
  \frac{3}{2\pi}\left(\frac{9\pi}{4}\right)^{1/3} = 0.9163,
  \label{eq:6-egcoefficients}
\end{equation}
as \texttt{hartreefock.py} confirms.  The first term is the kinetic energy,
which scales as $\rho^{2/3}\propto r_s^{-2}$ and dominates at high density;
the second is exchange, which scales as $\rho^{1/3}\propto r_s^{-1}$ and wins
at low density.  Because they carry opposite signs and different powers, their
sum has a minimum:
\begin{equation}
  \frac{\ud{}}{\ud{r_s}}\left[\frac{2.21}{r_s^2}
    - \frac{0.916}{r_s}\right] = 0
  \qquad\Longrightarrow\qquad
  r_s = \frac{2\times2.21}{0.916} = 4.83,
  \label{eq:6-egminimum}
\end{equation}
at which $E_0/N = -0.0949$~Ry $=-1.29$~eV.  The values along the curve are

\begin{center}
\renewcommand{\arraystretch}{1.2}
\begin{tabular}{rr}
\toprule
$r_s$ & $E_0/N$ [Ry]\\
\midrule
$1.00$ & $\phantom{-}1.294000$\\
$2.00$ & $\phantom{-}0.094500$\\
$3.00$ & $-0.059778$\\
$4.00$ & $-0.090875$\\
$4.83$ & $-0.094916$\\
$6.00$ & $-0.091278$\\
$8.00$ & $-0.079969$\\
\bottomrule
\end{tabular}
\end{center}

The electron gas is therefore \emph{bound} already at the Hartree-Fock level,
at a density lying squarely in the range observed in the alkali metals.  This
is a real success: exchange alone, with no correlation whatsoever, explains
why a metal is stable against expansion.  The binding energy is too small by
roughly a factor of two -- the missing part is the correlation energy, which
for this system is about $-0.06$~Ry near equilibrium -- and the equilibrium
density is somewhat too low, but the mechanism is right.

%----------------------------------------------------------------------
\section{Summary}
\label{sec:ch6summary}

Hartree-Fock theory is the variational principle of
Section~\ref{sec:variational} applied to the smallest interesting trial space:
a single Slater determinant, with the orbitals free.  The stationarity
condition is a nonlinear eigenvalue problem,
Eq.~(\ref{eq:6-hfequations}), solved by iterating a linear one -- the
self-consistent field.  Its solution is characterised by Brillouin's theorem,
$\bracket{\Phi_0}{\hat{H}}{\Phi_i^a}=0$, which is exactly the statement that
the reference decouples from all single excitations and which is why every
post-Hartree-Fock method begins at the doubles.

Thouless' theorem makes the geometry precise: the neighbourhood of a
determinant, inside the manifold of determinants, is parametrised by
$e^{\hat{T}}$ with $\hat{T}$ a linear combination of one-particle-one-hole
operators.  Because those operators commute, the exponential collapses, and
the Baker-Campbell-Hausdorff series -- which controls the general case and
reappears in the Trotter splitting and in coupled-cluster theory -- reduces
to its first term.  Going to second order in $\hat{T}$ produces the stability
matrix~(\ref{eq:6-stabilitymatrix}), whose positivity decides whether the
stationary point is a minimum, and which is the same matrix that the
random-phase approximation diagonalises.

The two model calculations bracket the range of behaviour.  For the pairing
model the mean field does nothing at all, since the interaction conserves
seniority.  For the Lipkin model it does too much: beyond $\chi=1$ the
symmetric solution goes unstable and a symmetry-broken one takes over, buying
energy at the cost of a broken symmetry that must later be restored.  Between
them lies the electron gas, where the whole programme goes through
analytically, predicts a bound system at a realistic density, and yet gets the
effective mass qualitatively wrong for want of screening.

The program \texttt{hartreefock.py} carries out every calculation of this
chapter, and the notebook \texttt{hartreefock.ipynb} runs it cell by cell.

%----------------------------------------------------------------------
\section{Exercises}
\label{sec:ch6exercises}
%=============================================================================

\subsection*{Exercise 1: determinants and unitary transformations}
\label{ex:6-determinants}
\addcontentsline{toc}{subsection}{Exercise 1: determinants and unitary transformations}

\begin{enumerate}
  \item[a)] Verify Eq.~(\ref{eq:6-multilinear2}) by direct expansion, and
  then prove Eq.~(\ref{eq:6-multilinear}) for a general $n$.
  \item[b)] Use a) to prove the product rule
  $\det(\bm{B}\bm{C})=\det(\bm{B})\det(\bm{C})$ for $n=2$, and convince
  yourself that the argument generalises.
  \item[c)] Show that a Slater determinant is invariant, up to a phase, under
  any unitary mixing of its occupied orbitals, and that the one-body density
  matrix~(\ref{eq:6-densitymatrix}) is strictly invariant.
  \item[d)] Show that $\bm{\rho}^2=\bm{\rho}$ and
  $\mathrm{tr}\,\bm{\rho}=N$, and explain what these two statements say about
  the occupied subspace.
\end{enumerate}

\paragraph{Answer to a).}
Expanding the left-hand side of Eq.~(\ref{eq:6-multilinear2}),
\begin{equation}
  (\alpha_1b_{11}+\alpha_2b_{12})a_{22}
  - (\alpha_1b_{21}+\alpha_2b_{22})a_{12}
  = \alpha_1(b_{11}a_{22}-b_{21}a_{12})
  + \alpha_2(b_{12}a_{22}-b_{22}a_{12}),
  \label{eq:6-exa1}
\end{equation}
which is the right-hand side.  For general $n$, use the Leibniz formula
$\det(\bm{A})=\sum_{\pi}\mathrm{sgn}(\pi)\prod_j a_{\pi(j)j}$.  Exactly one
factor in each product comes from column $j$; substituting
$a_{ij}\to\sum_k c_kb_{ik}$ in that column and interchanging the two sums
gives $\sum_k c_k\sum_{\pi}\mathrm{sgn}(\pi)b_{\pi(j)k}\prod_{l\ne j}
a_{\pi(l)l}$, which is Eq.~(\ref{eq:6-multilinear}).

\paragraph{Answer to b).}
Write $\bm{A}=\bm{B}\bm{C}$, so that column $j$ of $\bm{A}$ is
$\sum_k C_{kj}\bm{b}_k$ with $\bm{b}_k$ the columns of $\bm{B}$.  Applying
Eq.~(\ref{eq:6-multilinear}) to both columns of a $2\times2$ matrix,
\begin{equation}
  \det(\bm{A}) = \sum_{k_1k_2}C_{k_11}C_{k_22}
    \det\left(\bm{b}_{k_1}\ \bm{b}_{k_2}\right) .
  \label{eq:6-exb1}
\end{equation}
Terms with $k_1=k_2$ vanish, and the two surviving terms combine to
$(C_{11}C_{22}-C_{21}C_{12})\det(\bm{b}_1\ \bm{b}_2)
=\det(\bm{C})\det(\bm{B})$.  For general $n$ the same expansion leaves only
the $n!$ terms with distinct $k$'s, each carrying the sign of the
corresponding permutation, and the sum is again $\det(\bm{C})\det(\bm{B})$.

\paragraph{Answer to c).}
Let $\psi'_i=\sum_j U_{ij}\psi_j$ with $\bm{U}$ unitary and $i,j$ running over
the occupied orbitals.  By Eq.~(\ref{eq:6-detcdetphi}) the new determinant is
$\det(\bm{U})$ times the old one, and $|\det(\bm{U})|=1$ for a unitary matrix,
so the two states differ by a phase and describe the same physics.  The
density matrix is $\bm{\rho}'=\bm{U}^{\dagger}\bm{\rho}\bm{U}$ evaluated on
the occupied block; since $\bm{\rho}$ is the projector onto the span of the
occupied orbitals and $\bm{U}$ merely rotates within that span, $\bm{\rho}$ is
unchanged exactly, not merely up to a phase.

\paragraph{Answer to d).}
With orthonormal orbitals,
$\rho_{\gamma\delta}=\sum_i C_{i\gamma}C^{*}_{i\delta}$ gives
\begin{equation}
  (\bm{\rho}^2)_{\gamma\delta}
  = \sum_{i j\lambda}C_{i\gamma}C^{*}_{i\lambda}C_{j\lambda}C^{*}_{j\delta}
  = \sum_{ij}\delta_{ij}C_{i\gamma}C^{*}_{j\delta}
  = \rho_{\gamma\delta},
  \label{eq:6-exd1}
\end{equation}
and $\mathrm{tr}\,\bm{\rho}=\sum_i\sum_{\gamma}|C_{i\gamma}|^2=N$.  A
Hermitian idempotent is a projector, and its trace is the dimension of the
space it projects onto.  So $\bm{\rho}$ is the projector onto the
$N$-dimensional occupied subspace, and its eigenvalues are $N$ ones and the
rest zeros.  Together these say that a Hartree-Fock solution is a choice of
subspace, and that only for a Slater determinant are the occupation numbers
exactly zero and one -- for a correlated state the eigenvalues of $\bm{\rho}$
are the fractional natural occupations of
Section~\ref{sec:naturalorbitals}.

\subsection*{Exercise 2: the Hartree-Fock equations}
\label{ex:6-hfequations}
\addcontentsline{toc}{subsection}{Exercise 2: the Hartree-Fock equations}

\begin{enumerate}
  \item[a)] Carry out the differentiation in Eq.~(\ref{eq:6-variationC})
  in full, and confirm that the factor $\tfrac12$ in the two-body term is
  cancelled by the two ways in which $C^{*}_{i\alpha}$ can appear.
  \item[b)] Show that the Hartree-Fock energy is \emph{not} the sum of the
  single-particle energies, and derive the correct relation
  \begin{equation*}
    E^{\mathrm{HF}} = \sum_{i=1}^{N}\varepsilon_i^{\mathrm{HF}}
      - \frac{1}{2}\sum_{i,j=1}^{N}
        \langle ij\vert\hat{v}\vert ij\rangle_{\mathrm{AS}} .
  \end{equation*}
  \item[c)] Interpret $\varepsilon_i^{\mathrm{HF}}$ physically by computing
  the energy needed to remove a particle from orbital $i$, assuming the
  remaining orbitals do not relax.  (This is Koopmans' theorem.)
  \item[d)] Show directly from Eq.~(\ref{eq:6-fockmatrix}) that at
  self-consistency $f_{ai}=0$, and hence that
  $\bracket{\Phi_0}{\hat{H}}{\Phi_i^a}=0$.
\end{enumerate}

\paragraph{Answer to a).}
The one-body term of Eq.~(\ref{eq:6-Efunctional}) is linear in
$C^{*}_{i\alpha}$ and differentiates to
$\sum_{\beta}C_{i\beta}\bracket{\alpha}{\hat{h}_0}{\beta}$.  In the two-body
term $C^{*}_{i\alpha}$ can be either the first starred factor, with the sum
over $j$ free, or the second, with the roles of the index pairs exchanged.
Relabelling the dummy indices of the second contribution and using
\begin{equation}
  \langle\alpha\beta\vert\hat{v}\vert\gamma\delta\rangle_{\mathrm{AS}}
  = \langle\beta\alpha\vert\hat{v}\vert\delta\gamma\rangle_{\mathrm{AS}}
  \label{eq:6-exa2}
\end{equation}
shows the two are identical, so together they cancel the $\tfrac12$ and give
the second term of Eq.~(\ref{eq:6-hfraw}).

\paragraph{Answer to b).}
Multiply Eq.~(\ref{eq:6-hfequations}) by $C^{*}_{i\alpha}$ and sum over
$\alpha$ and over the occupied $i$:
\begin{equation}
  \sum_{i=1}^{N}\varepsilon_i^{\mathrm{HF}}
  = \sum_{i=1}^{N}\bracket{i}{\hat{h}_0}{i}
  + \sum_{i,j=1}^{N}\langle ij\vert\hat{v}\vert ij\rangle_{\mathrm{AS}} .
  \label{eq:6-exb2}
\end{equation}
Comparing with Eq.~(\ref{eq:6-energyfunctional}), which carries a factor
$\tfrac12$ on the two-body term, gives the stated relation.  The interaction
energy is counted twice in the sum of single-particle energies -- once for
each partner -- and half of it must be removed.  This is the standard
double-counting correction, and forgetting it is one of the commonest errors
in a first Hartree-Fock code.

\paragraph{Answer to c).}
Removing a particle from orbital $i$ without relaxing the others changes the
energy by
\begin{equation}
  E^{\mathrm{HF}}(N) - E^{\mathrm{HF}}(N-1)
  = \bracket{i}{\hat{h}_0}{i}
    + \sum_{j\ne i}\langle ij\vert\hat{v}\vert ij\rangle_{\mathrm{AS}}
  = \varepsilon_i^{\mathrm{HF}},
  \label{eq:6-exc2}
\end{equation}
using $\langle ii\vert\hat{v}\vert ii\rangle_{\mathrm{AS}}=0$.  So each
Hartree-Fock eigenvalue is, in the frozen-orbital approximation, minus the
ionisation energy from that orbital.  This is Koopmans' theorem.  It is only
an approximation, because in reality the remaining orbitals relax -- an effect
that lowers the energy of the $(N-1)$-particle system and therefore makes the
true ionisation energy smaller than $-\varepsilon_i^{\mathrm{HF}}$.
Correlation works in the opposite direction, and for atoms the two errors
partly cancel, which is why Koopmans' theorem is more accurate than it has any
right to be.

\paragraph{Answer to d).}
At self-consistency the coefficients diagonalise the Fock matrix, so in the
Hartree-Fock basis $f_{pq}=\varepsilon_p\delta_{pq}$ and in particular
$f_{ai}=0$ for $a>F\ge i$.  By the Condon-Slater rules of
Section~\ref{sec:wickapps}, the matrix element of $\hat{H}$ between the
reference and a single excitation is
\begin{equation}
  \bracket{\Phi_0}{\hat{H}}{\Phi_i^a}
  = \bracket{a}{\hat{h}_0}{i}
    + \sum_{j\le F}\langle aj\vert\hat{v}\vert ij\rangle_{\mathrm{AS}}
  = f_{ai} = 0 .
  \label{eq:6-exd2}
\end{equation}
This is the identity that makes the Hartree-Fock determinant the natural
reference for everything that follows.

\subsection*{Exercise 3: Thouless' theorem}
\label{ex:6-thouless}
\addcontentsline{toc}{subsection}{Exercise 3: Thouless' theorem}

\begin{enumerate}
  \item[a)] Prove Eq.~(\ref{eq:6-phcommute}) from the anticommutation
  relations, being explicit about where the disjointness of the particle and
  hole index sets is used.
  \item[b)] Show that $\hat{T}_i^2=0$ and hence that
  $e^{\hat{T}}=\prod_i(1+\hat{T}_i)$.  Explain why this is \emph{not} the
  same as $1+\hat{T}$, and identify what the difference describes.
  \item[c)] Verify all of this numerically with \texttt{hartreefock.py} for
  six orbitals and three particles, and check that the discarded piece scales
  as the square of the amplitudes.
  \item[d)] Show that $\braket{c}{\tilde{c}}=\det(f_{ip})$ over the occupied
  block, and explain why the hypothesis $\braket{c}{\tilde{c}}\ne0$ is exactly
  what is needed for the proof to go through.
\end{enumerate}

\paragraph{Answer to a).}
Take $a,b>F$ and $i,j\le F$, so that all four indices in a given product are
distinct except possibly $a=b$ or $i=j$.  Since $i\ne b$,
$\hat{a}_{i}\hat{a}^{\dagger}_{b}=-\hat{a}^{\dagger}_{b}\hat{a}_{i}$, whence
\begin{equation}
  \hat{a}^{\dagger}_{a}\hat{a}_{i}\hat{a}^{\dagger}_{b}\hat{a}_{j}
  = -\hat{a}^{\dagger}_{a}\hat{a}^{\dagger}_{b}\hat{a}_{i}\hat{a}_{j} .
  \label{eq:6-exa3}
\end{equation}
Since $j\ne a$, the same step on the reversed product gives
\begin{equation}
  \hat{a}^{\dagger}_{b}\hat{a}_{j}\hat{a}^{\dagger}_{a}\hat{a}_{i}
  = -\hat{a}^{\dagger}_{b}\hat{a}^{\dagger}_{a}\hat{a}_{j}\hat{a}_{i}
  = -\hat{a}^{\dagger}_{a}\hat{a}^{\dagger}_{b}\hat{a}_{i}\hat{a}_{j},
  \label{eq:6-exb3}
\end{equation}
where the second equality interchanges both the two creation operators and
the two annihilation operators, contributing two signs that cancel.  The two
products are equal, so the commutator vanishes.  The disjointness of the index
sets is what allowed both anticommutations to be performed without generating
a delta function; this is why the result is special to particle-hole
operators and fails for general $\hat{a}^{\dagger}_p\hat{a}_q$.

\paragraph{Answer to b).}
$\hat{T}_i^2$ contains $\hat{a}_{i}\cdots\hat{a}_{i}$, and by
Eq.~(\ref{eq:6-exa3}) the two annihilation operators can be brought adjacent;
$\hat{a}_i^2=0$, so $\hat{T}_i^2=0$ and $e^{\hat{T}_i}=1+\hat{T}_i$.
Combining with a) gives
$e^{\hat{T}}=\prod_i e^{\hat{T}_i}=\prod_i(1+\hat{T}_i)$.  Expanding the
product,
\begin{equation}
  \prod_i(1+\hat{T}_i) = 1 + \sum_i\hat{T}_i
    + \sum_{i<j}\hat{T}_i\hat{T}_j + \cdots
    = 1 + \hat{T} + \bigO(\hat{T}^2),
  \label{eq:6-exb4}
\end{equation}
so the difference from $1+\hat{T}$ consists of the cross terms
$\hat{T}_i\hat{T}_j$ with $i\ne j$, which are $2p$-$2h$, $3p$-$3h$, \dots\
excitations.  They are second and higher order in the amplitudes, which is
why the linearisation is exact for an infinitesimal variation and why the
stability analysis of Section~\ref{sec:hfstability}, which works to second
order in the energy but first order in the state, is consistent.

\paragraph{Answer to c).}
The class \texttt{ThoulessRotation} reports
$\max|[\hat{T}_i,\hat{T}_j]|=2.8\times10^{-17}$,
$\max|\hat{T}_i^2|=0$ exactly, agreement between $e^{\hat{T}}\bket{c}$ and
$\prod_i\hat{b}^{\dagger}_i\bket{0}$ to $5.6\times10^{-17}$, an overlap
$\braket{c}{c'}=1$ and $\det(f)=1$ over the occupied block.  Scaling the
amplitudes by $s$ and comparing $e^{s\hat{T}}\bket{c}$ with
$(1+s\hat{T})\bket{c}$ gives errors $3.68\times10^{-2}$,
$9.20\times10^{-3}$, $2.30\times10^{-3}$, $5.75\times10^{-4}$ and
$1.44\times10^{-4}$ for $s=0.4,0.2,0.1,0.05,0.025$: each halving of $s$
divides the error by four, confirming the $\bigO(s^2)$ scaling.

\paragraph{Answer to d).}
Expanding Eq.~(\ref{eq:6-generalb}) and taking the overlap with
$\bbra{c}=\bbra{0}\hat{a}_{i_n}\cdots\hat{a}_{i_1}$, only terms in which every
$\hat{a}^{\dagger}_p$ carries an index below the Fermi level survive, and each
must be contracted with a distinct $\hat{a}_{i_k}$.  Summing over the
assignments with their signs reconstructs
$\sum_{\pi}\mathrm{sgn}(\pi)\prod_k f_{i_k\pi(i_k)}=\det(f_{ip})$ over the
occupied block.  The proof needs $f$ restricted to that block to be
invertible, since $f^{-1}$ appears in Eq.~(\ref{eq:6-finverse}); a
non-vanishing determinant is precisely invertibility.  Geometrically, if
$\braket{c}{\tilde{c}}=0$ the two determinants have no overlap and
$\bket{\tilde{c}}$ is not in the neighbourhood of $\bket{c}$ at all -- one
cannot reach it by an infinitesimal rotation, and Thouless' theorem makes no
claim about it.

\subsection*{Exercise 4: BCH and Trotter}
\label{ex:6-bch}
\addcontentsline{toc}{subsection}{Exercise 4: BCH and Trotter}

\begin{enumerate}
  \item[a)] Derive Eq.~(\ref{eq:6-bchA2}) by matching the power series of
  $e^{X}e^{Y}$ and $e^{Z}$ to second order.
  \item[b)] Show that if $[X,Y]=cI$ then Eq.~(\ref{eq:6-bchcentral}) is
  exact, and use it to derive the Weyl form
  $e^{\imath a\hat{x}/\hbar}e^{\imath b\hat{p}/\hbar}
  = e^{-\imath ab/2\hbar}e^{\imath(a\hat{x}+b\hat{p})/\hbar}$.
  \item[c)] Use Eq.~(\ref{eq:6-bch}) to show that the leading error of the
  first-order splitting is $-\tfrac{1}{2}[H_1,H_2]\Delta^2$, and verify that
  the symmetric formula~(\ref{eq:6-strang}) cancels it.  Which commutators
  govern the leading error of $S_2$?
  \item[d)] Prove the Baker-Hausdorff lemma
  $e^{-X}Ye^{X}=Y-[X,Y]+\tfrac1{2!}[X,[X,Y]]-\cdots$ and explain why the
  corresponding series in coupled-cluster theory terminates.
  \item[e)] Reproduce the two tables of Sections~\ref{sec:bch}
  and~\ref{sec:trotter} with \texttt{hartreefock.py}, and estimate how many
  Trotter steps a first- and a second-order scheme need for an accuracy
  $\varepsilon=10^{-3}$ at $t=1$.
\end{enumerate}

\paragraph{Answer to a).}
See Eq.~(\ref{eq:6-bchmatch}) and the discussion following it.  The essential
observation is that $\tfrac12(X+Y)^2$ symmetrises the cross term into
$\tfrac12(XY+YX)$, whereas $e^{X}e^{Y}$ produces $XY$ with unit weight; the
deficit is $\tfrac12(XY-YX)$.

\paragraph{Answer to b).}
If $[X,Y]=cI$ then $[X,[X,Y]]=[Y,[X,Y]]=0$ and every term of
Eq.~(\ref{eq:6-bch}) beyond $\tfrac12[X,Y]$ contains at least one such double
commutator.  For $X=\imath a\hat{x}/\hbar$ and $Y=\imath b\hat{p}/\hbar$ we
have $[X,Y]=-(ab/\hbar^2)[\hat{x},\hat{p}]=-\imath ab/\hbar$, so
\begin{equation}
  e^{X}e^{Y} = \exp\left(\frac{\imath}{\hbar}(a\hat{x}+b\hat{p})
    - \frac{\imath ab}{2\hbar}\right)
  = e^{-\imath ab/2\hbar}\,
    e^{\imath(a\hat{x}+b\hat{p})/\hbar}.
  \label{eq:6-exb5}
\end{equation}
The same argument applied to $[\hat{a},\hat{a}^{\dagger}]=1$ gives the
normal-ordered form of the coherent-state displacement operator.

\paragraph{Answer to c).}
Setting $X=-\imath H_1\Delta$, $Y=-\imath H_2\Delta$ in
Eq.~(\ref{eq:6-bch}), the second term is
$\tfrac12[X,Y]=-\tfrac12[H_1,H_2]\Delta^2$, which is
Eq.~(\ref{eq:6-trotterbch}).  For the symmetric product, write
$S_2(\Delta)=e^{X/2}e^{Y}e^{X/2}$; applying BCH twice, the $[X,Y]$
contributions from the two half-steps enter with opposite signs and cancel
exactly.  What survives at the next order is a combination of
$[H_1,[H_1,H_2]]$ and $[H_2,[H_1,H_2]]$ multiplying $\Delta^3$.  This is also
visible from symmetry alone: $S_2(-\Delta)=S_2(\Delta)^{-1}$, so the error
must be an odd function of $\Delta$ and cannot contain a $\Delta^2$ term.

\paragraph{Answer to d).}
Define $G(\lambda)=e^{-\lambda X}Ye^{\lambda X}$.  Then
$G'(\lambda)=e^{-\lambda X}[Y,X]e^{\lambda X}=-e^{-\lambda X}[X,Y]
e^{\lambda X}$, and iterating,
$G^{(n)}(0)=(-1)^n\,\mathrm{ad}_X^n(Y)$ with
$\mathrm{ad}_X(Y)=[X,Y]$.  Taylor-expanding about $\lambda=0$ and setting
$\lambda=1$ gives the lemma.  In coupled-cluster theory $X=\hat{T}$ is a
sum of excitation operators, all of which commute with one another, and
$Y=\hat{H}$ contains at most two-body terms.  Each commutator with $\hat{T}$
must contract at least one annihilation operator of $\hat{H}$, and $\hat{H}$
has only four, so the series terminates exactly after the fourth commutator --
which is why coupled cluster is a finite, not an asymptotic, expansion.

\paragraph{Answer to e).}
The tables are reproduced by \texttt{demo\_bch()} and
\texttt{demo\_trotter()}.  Reading off the first-order column, the error is
about $1.0/N$, so $\varepsilon=10^{-3}$ requires $N\approx10^{3}$ steps.  The
second-order error is about $0.39/N^2$, giving $N\approx20$.  The second-order
scheme costs about $50\%$ more exponentials per step, so it is cheaper by more
than an order of magnitude here -- and the advantage grows as the accuracy
demand tightens, which is why practical quantum-simulation circuits are never
built at first order.

\subsection*{Exercise 5: stability and the Lipkin model}
\label{ex:6-stability}
\addcontentsline{toc}{subsection}{Exercise 5: stability and the Lipkin model}

\begin{enumerate}
  \item[a)] Derive Eq.~(\ref{eq:6-cprimenorm}) for the norm of the rotated
  determinant, keeping track of which terms are second order.
  \item[b)] Starting from Eq.~(\ref{eq:6-energyquadratic}), obtain
  Eq.~(\ref{eq:6-deltaE}) and the stability
  matrix~(\ref{eq:6-stabilitymatrix}).  Why does $\chi$ contain both
  $\delta C$ and $\delta C^{*}$?
  \item[c)] For the Lipkin model with $W=0$, show that $A=0$ and that the
  largest eigenvalue of $B$ is $|V|(N-1)$, and hence recover
  Eq.~(\ref{eq:6-lipkinlowest}).
  \item[d)] Derive the mean-field energy~(\ref{eq:6-lipkinmeanfield}) from
  the rotation~(\ref{eq:6-lipkinrotation}), and locate both stationary
  points.
  \item[e)] Reproduce Table~\ref{tab:lipkinstability} and comment on the gap
  between $E_{\mathrm{mf}}$ and $E_{\mathrm{exact}}$ on either side of
  $\chi=1$.
\end{enumerate}

\paragraph{Answer to a).}
Expanding Eq.~(\ref{eq:6-cprimeexpansion}) and taking the overlap with
itself, the zeroth-order term gives $1$; the cross terms between the constant
and the linear piece vanish, since
$\bracket{c}{\hat{a}^{\dagger}_a\hat{a}_i}{c}=0$ for $a>F\ge i$; and the
linear piece with itself gives
$\sum_{abij}\delta C^{*}_{ai}\delta C_{bj}\braket{\Phi_i^a}{\Phi_j^b}
=\sum_{ai}|\delta C_{ai}|^2$ by the orthonormality of the excited
determinants.  Terms involving the quadratic piece of the expansion are of
third order or higher when paired with the linear one, and of fourth order
with themselves.

\paragraph{Answer to b).}
Subtract $(1+\sum|\delta C|^2)E_0$ from Eq.~(\ref{eq:6-energyquadratic}) and
divide by the norm; what is left over the norm is $\Delta E$.  Collecting the
terms, those proportional to $\delta C^{*}\delta C$ assemble into
$\chi^{\dagger}$ times the diagonal blocks $\Delta+A$, and those proportional
to $\delta C\delta C$ and $\delta C^{*}\delta C^{*}$ into the off-diagonal
blocks $B$ and $B^{*}$.  The doubled vector $\chi$ is needed precisely because
the quadratic form is not Hermitian in $\delta C$ alone: the terms
$B\,\delta C^{*}\delta C^{*}$ pair two conjugated amplitudes, and only by
treating $\delta C$ and $\delta C^{*}$ as independent components of a single
vector can the whole expression be written as a Hermitian form.  This doubling
is exactly the structure of the RPA equations, where it produces the pairs of
solutions $\pm\omega$.

\paragraph{Answer to c).}
With $W=0$ the interaction is $(V/2)(\hat{J}_+^2+\hat{J}_-^2)$, which changes
$J_z$ by $\pm2$, so it has no matrix element within the $1p$-$1h$ space and
$A_{ai,bj}=0$; the diagonal block is $\Delta=\varepsilon\bm{I}$, since every
particle-hole excitation costs $\varepsilon$.  For $B$, the operator
$\hat{J}_-^2$ connects the reference to a $2p$-$2h$ state in which levels $p$
and $q$ are both excited, giving $B_{ai,bj}=V$ for $p\ne q$ and zero for
$p=q$; that is, $\bm{B}=V(\bm{J}-\bm{I})$ within each spin block, with
$\bm{J}$ the matrix of ones.  Its eigenvalues are $V(N-1)$, once, and $-V$,
$N-1$ times.  Because $\hat{M}$ has the form
$\left(\begin{smallmatrix}\varepsilon\bm{I} & \bm{B}\\ \bm{B} &
\varepsilon\bm{I}\end{smallmatrix}\right)$, its eigenvalues are
$\varepsilon\pm b$ with $b$ an eigenvalue of $\bm{B}$, and the smallest is
$\varepsilon-|V|(N-1)=\varepsilon(1-\chi)$.

\paragraph{Answer to d).}
Under the rotation~(\ref{eq:6-lipkinrotation}) the quasispin vector rotates
rigidly, $\langle\hat{J}_z\rangle=-\tfrac{N}{2}\cos\alpha$ and
$\langle\hat{J}_x\rangle=\tfrac{N}{2}\sin\alpha$.  The one-body term gives
$-\tfrac{N\varepsilon}{2}\cos\alpha$ directly.  For the interaction,
$\hat{J}_+^2+\hat{J}_-^2=2(\hat{J}_x^2-\hat{J}_y^2)$, and evaluating this in
the rotated determinant -- where the $N$ particles are uncorrelated, so that
$\langle\hat{J}_x^2\rangle=\langle\hat{J}_x\rangle^2+\bigO(N)$ -- gives
$-\tfrac{N\varepsilon\chi}{4}\sin^2\alpha$ with
$\chi=|V|(N-1)/\varepsilon$, the factor $N-1$ rather than $N$ arising because
a particle does not interact with itself.  Setting the
derivative~(\ref{eq:6-lipkinderivative}) to zero gives $\sin\alpha=0$ or
$\cos\alpha=1/\chi$, the second having a solution only when $\chi\ge1$.

\paragraph{Answer to e).}
The table is produced by \texttt{demo\_stability()}.  Below $\chi=1$ the
mean-field energy is flat at $-N\varepsilon/2=-2$ while the exact energy falls
steadily, so the error grows: the mean field is stable but increasingly
inadequate, and the missing energy is pure correlation.  At $\chi=1.5$ and
$2.0$ the deformed solution recovers part of it -- $0.167$ and $0.5$
respectively -- but remains $0.48$ and $0.56$ above the exact answer.  The
deformed state breaks the symmetry of $\hat{H}$ under the rotation generated
by $\hat{J}_z$ parity, which the exact ground state respects, so part of the
remaining error is spurious and would be removed by projecting the broken
symmetry back out.

\subsection*{Exercise 6: the electron gas}
\label{ex:6-electrongas}
\addcontentsline{toc}{subsection}{Exercise 6: the electron gas}

\begin{enumerate}
  \item[a)] Derive Eq.~(\ref{eq:6-Vmomentum}) for a general interaction
  $v(r)$, and specialise it to the screened Coulomb interaction.
  \item[b)] Show that $\hat{H}_b$ and $\hat{H}_{\mathrm{el}-b}$ take the
  forms~(\ref{eq:6-divergences}), and that together with the direct term of
  the electron-electron repulsion they cancel in the thermodynamic limit.
  \item[c)] Carry out the steps from Eq.~(\ref{eq:6-egsp}) to
  Eq.~(\ref{eq:6-egspclosed}).
  \item[d)] Compute the Hartree-Fock effective mass $m^{*}_{\mathrm{HF}}$ and
  the level density $n(\varepsilon_F^{\mathrm{HF}})$, and comment on the
  result.
  \item[e)] Show that the energy per electron is
  Eq.~(\ref{eq:6-egenergy}), plot it, and explain why the system is bound.
  Compute the pressure $P=-(\partial E/\partial\Omega)_N$ and the bulk
  modulus $B=-\Omega(\partial P/\partial\Omega)_N$, and comment.
\end{enumerate}

\paragraph{Answer to a).}
See Eqs.~(\ref{eq:6-planewaveme})--(\ref{eq:6-Vmomentum}).  For the screened
Coulomb interaction the Fourier transform is
\begin{equation}
  \int \frac{e^2e^{-\mu r}}{r}e^{-\imath\bm{q}\cdot\bm{r}}\ud{\bm{r}}
  = \frac{4\pi e^2}{\mu^2+q^2}
  \ \xrightarrow[\mu\to0]{}\ \frac{4\pi e^2}{q^2},
  \label{eq:6-exa6}
\end{equation}
which is Eq.~(\ref{eq:6-egstep2}) again.  The $q^{-2}$ divergence at small
momentum transfer is the signature of the infinite range of the interaction,
and it is the source of every difficulty in this section.

\paragraph{Answer to b).}
Both are elementary integrals of $e^{-\mu r}/r$ over the box with a constant
density $n=N/\Omega$: $\int \ud{\bm{r}}\,e^{-\mu r}/r=4\pi/\mu^2$, whence
$\hat{H}_b=\tfrac{e^2}{2}n^2\Omega\cdot 4\pi/\mu^2
=\tfrac{e^2}{2}\tfrac{N^2}{\Omega}\tfrac{4\pi}{\mu^2}$ and
$\hat{H}_{\mathrm{el}-b}=-e^2\tfrac{N^2}{\Omega}\tfrac{4\pi}{\mu^2}$, twice as
large and of opposite sign.  The direct term of
Eq.~(\ref{eq:6-egpotential}) is $\tfrac{N^2}{2\Omega}\int v(r)\ud{\bm{r}}
=\tfrac{e^2}{2}\tfrac{N^2}{\Omega}\tfrac{4\pi}{\mu^2}$, and the three sum to
zero.  Physically: a neutral system has no monopole, so the $\bm{q}=0$
component of the interaction must drop out, and it is only the separation into
three individually divergent pieces that made the cancellation look accidental.

\paragraph{Answer to c).}
See Eqs.~(\ref{eq:6-egstep1})--(\ref{eq:6-egstep4}).  The last integral is
\begin{equation}
  \int_0^{k_F}p\ln\left\vert\frac{p+k}{p-k}\right\vert\ud{p}
  = kk_F + \frac{k_F^2-k^2}{2}
    \ln\left\vert\frac{k+k_F}{k-k_F}\right\vert ,
  \label{eq:6-exc6}
\end{equation}
which may be checked by differentiating with respect to $k_F$, or numerically.

\paragraph{Answer to d).}
Differentiating Eq.~(\ref{eq:6-egdimensionless}),
\begin{equation}
  \frac{m^{*}_{\mathrm{HF}}}{m}
  = \frac{2}{2-0.663\,r_s\,F'(1)},
  \qquad
  F'(x) \sim -\frac{1}{2}\ln\left\vert\frac{2}{1-x}\right\vert
  \ \ \text{as } x\to1 .
  \label{eq:6-exd6}
\end{equation}
Since $F'(1)=-\infty$, the denominator diverges and
$m^{*}_{\mathrm{HF}}\to0$; the program gives $0.256$, $0.185$, $0.144$,
$0.118$ and $0.100$ at $1-x=10^{-2},\dots,10^{-6}$, decreasing like the
reciprocal logarithm.  By Eq.~(\ref{eq:6-leveldensity}) the level density is
proportional to $m^{*}$ and therefore vanishes at the Fermi surface.  This
contradicts experiment: the electronic specific heat of a metal is linear in
$T$ with a coefficient proportional to $n(\varepsilon_F)$, and it is measured
to be finite.  The failure is traceable to the unscreened $1/q^2$, and it is
repaired by the random-phase approximation, which replaces it by a screened
interaction with a finite range.

\paragraph{Answer to e).}
The kinetic term is Eq.~(\ref{eq:6-egkineticdensity}); writing
$k_F=(9\pi/4)^{1/3}/r_0$ and measuring in Rydberg,
\begin{equation}
  \frac{\langle\hat{T}\rangle}{N}
  = \frac{3}{5}\frac{\hbar^2k_F^2}{2m}
  = \frac{e^2}{2a_0}\,\frac{3}{5}
    \left(\frac{9\pi}{4}\right)^{2/3}\frac{1}{r_s^2}
  = \frac{e^2}{2a_0}\frac{2.21}{r_s^2},
  \label{eq:6-exe6}
\end{equation}
and the exchange term evaluates similarly to $-0.916/r_s$ in the same units.
The system is bound because the two terms carry opposite signs and different
powers of $r_s$: the kinetic repulsion, $\propto r_s^{-2}$, dominates at high
density and the exchange attraction, $\propto r_s^{-1}$, at low density, so
their sum has a minimum, at $r_s=4.83$ with $E_0/N=-0.0949$~Ry.  For the
pressure, with $\Omega\propto Nr_0^3$ so that
$\partial/\partial\Omega=-(r_s/3\Omega)\,\partial/\partial r_s$,
\begin{equation}
  P = -\left(\frac{\partial E}{\partial\Omega}\right)_N
    = \frac{\rho\,r_s}{3}\frac{\ud{}}{\ud{r_s}}
      \left(\frac{E_0}{N}\right)
    = \frac{e^2\rho}{2a_0}\frac{1}{3}
      \left[-\frac{4.42}{r_s^2}+\frac{0.916}{r_s}\right],
  \label{eq:6-exe7}
\end{equation}
which vanishes exactly at the equilibrium $r_s$, as it must, is negative
(the gas would contract) at higher density and positive at lower.  The bulk
modulus $B=-\Omega(\partial P/\partial\Omega)_N$ is positive at the minimum,
confirming mechanical stability, and its measured values in the alkali metals
are of the right order but consistently smaller than the Hartree-Fock
prediction -- correlation softens the gas.
